\documentclass[11pt]{memoir}
\usepackage[usenames,dvipsnames]{color}
\usepackage{amsmath}%
\usepackage{amsfonts}%
\usepackage{amssymb}%
\usepackage{amsthm}%
\usepackage{graphicx}
\usepackage{dsfont}
\usepackage[hidelinks]{hyperref} %hyperlink
\usepackage{subcaption} %NON SO SE SERVE
\usepackage{enumitem} %enumerazioni con le i
\usepackage{booktabs} %tabelle più belle
\usepackage{placeins} %forza le immagini in posizione
\usepackage{mathabx}
\usepackage{epsfig}
\usepackage{caption} %caption fuori da ambienti float
\usepackage{appendix} %appendice per i codici
\usepackage{dsfont}
\usepackage{color}
\usepackage{listings}
\usepackage{xcolor} %per colorare il testo
\def\R{\mathbb{R}}
\def\Z{\mathbb{Z}}
\def\H{\mathbb{H}}
\def\A{\mathbb{A}d\mathbb{S}}
\def\D{\mathbb{D}}
\def\S{\mathbb{S}}
\def\AS{\tilde{\mathbb{A}}d\mathbb{S}}
\def\1{\mathds{1}}
\newcommand\norm[1]{\lVert#1\rVert}


\definecolor{mygreen}{RGB}{28,172,0} % color values Red, Green, Blue
\definecolor{mylilas}{RGB}{170,55,241}
\makeatletter
\makeatother


\usepackage[utf8]{inputenc}
\usepackage[english]{babel}
\usepackage{mathrsfs}
\usepackage{array}
\usepackage{rotating}
\usepackage{multirow}
\usepackage{amsmath,amsthm}
\usepackage{mathtools}
\usepackage{stmaryrd}
\usepackage{tikz-cd}
\usepackage{csquotes}
\usepackage{lscape}
\tikzcdset{
    startcode/.style={
        start anchor={[xshift=0.35cm]south west}
        },
    endcode/.style={
        end anchor={[xshift=0.35cm]north west}
    }
}
\usepackage{tikz}
\usepackage{upgreek}
\usepackage{varwidth}
\usepackage{hyperref}
\usepackage{biblatex}
\usepackage[colorinlistoftodos]{todonotes}
\addbibresource{refs.bib}
\usetikzlibrary{arrows}
\newcommand{\encadrer}[1]{{%
  
  \fbox{%
    \begin{varwidth}{\dimexpr\linewidth-2\fboxsep-2\fboxrule}
      #1%
    \end{varwidth}%
  }
  \par}%
}

\theoremstyle{plain}
\newtheorem{theorem}{Theorem}[chapter] % reset theorem numbering for each chapter
\theoremstyle{definition}
\newtheorem{definition}[theorem]{Definition} % definition numbers are dependent on theorem numbers
\newtheorem{lemma}[theorem]{Lemma} % same for example numbers
\theoremstyle{plain}
\newtheorem{proposition}[theorem]{Proposition}
\newtheorem{observation}[theorem]{Remark}
\renewcommand*{\proofname}{Proof}
\newcommand\restr[2]{\ensuremath{\left.#1\right|_{#2}}}
\newtheorem{corollary}[theorem]{Corollary}

\chapterstyle{veelo}
\checkandfixthelayout
\setcounter{tocdepth}{2}


\begin{document}
\advance\vsize by 2.5cm % Advance page height
\advance\voffset by -1.5cm % Shift top margin


\begin{titlingpage}


\advance\hsize by 2.7cm % Advance page height
\advance\hoffset by -1.35cm % Shift top margin


%%\aliaspagestyle{titlingpage}{empty}

%\newgeometry{margin=2.5cm}

\calccentering{\unitlength}
\begin{adjustwidth*}{\unitlength}{-\unitlength}


\thispagestyle{empty}
\begin{center}
\large
\textsc{\Large Università di Pisa\\}
\vspace{0.6cm}
\includegraphics[width=8cm]{cherubino.eps}

\vspace{0.8cm}
\textsc{\Large{Corso di laurea in Matematica}}\\



\vspace{2.5cm}


{\LARGE\textbf{Thurston Earthquake Theorem via Anti-de Sitter Geometry}}
\\[2.0cm]

\textsc{Tesi di Laurea Magistrale \\[0.2cm] in Matematica}\\
\vspace{0.5cm}

\begin{center}

\textsc{\Large Anno Accademico 2023/2024}
\end{center}

\vspace{0.2cm}

%\begin{center}
%\textsc{{\LARGE Candidato}}\\
%\textbf{{\small Tizio Caio}}
%\end{center}
%
%\vspace{0.4cm}
%
%\begin{center}
%\textsc{{\LARGE Relatore}}\\
%\textbf{{\small Sempronio Mevio}} \\
%{\small Università di Qualche Posto}
%\end{center}

\begin{center}
\makebox[\textwidth]{
\begin{tabular}{l p{5.3cm} r}
\textsc{{\LARGE Candidato}} & & \textsc{{\LARGE Relatore}} \\
\textbf{{\small SMC}} & & \textbf{{\small AT}} \\
  & & {\small Università di Pisa} \\
\end{tabular}}
\end{center}




\vspace{3.75cm}


\textsc{Anno Accademico 2023 - 2024}\\[0.2cm]


\end{center}

\end{adjustwidth*}


\advance\hsize by -2.7cm % Advance page height
\advance\hoffset by 1.35cm % Shift top margin


\end{titlingpage}


\advance\vsize by -2.5cm % Return old margings and page height
\advance\voffset by 1.5cm % Return old margings and page height



%\par\vfill\break % Break the page with different margins



%\restoregeometry

\par\vfill\break

\thispagestyle{empty}
\vspace*{\stretch{1}}
\begin{flushright}


\end{flushright}
\vspace{\stretch{2}}
\cleardoublepage

\clearpage

\clearpage

\tableofcontents
\clearpage

\thispagestyle{empty}
\listoffigures

\thispagestyle{empty}
%\include{legenda}

%\include{ringraziamenti}
\setcounter{chapter}{-1}
\chapter{Introduction}

As in Riemannian geometry the \textit{most interesting} geometry of constant sectional curvature is the hyperbolic one, in Lorentzian geometry the analogous (as in with constant negative curvature) is the \textit{Anti-de Sitter} geometry. In the first chapter we will introduce the basic concepts of Lorentzian geometry and explore the relationships between general spaces with constant sectional curvature and spaces with maximal isometry group and constant sectional curvature.
We would then like to expand the study of Anti-de Sitter geometry in chapter two with a brief exposition of the most common models of AdS spaces and the study of geodesics in this manifolds. 
In chapter three we would like to specialize to $2+1$ Anti-de Sitter geometry fixing the dimension and working with explicit models, the main result of this section is the identification of AdS space with groups of matrices. 
In chapter four, with this "newly introduced" concepts, we would like to recover a celebrated theorem of Thurston related to hyperbolic surfaces following the pioneeristic work of G. Mess \cite{Mess}.
\chapter{Lorentzian geometry}

In what follows we will introduce Lorentzian manifold of constant curvature and shows that, as in the Riemannian case, two manifolds of constant sectional curvature $K$ are locally isometric. We will then define what we mean by a manifold with maximal isometry group as they provide models of manifold with constant curvature: if we have $M$ a Lorentzian manifold with costant sectional curvature $K$ and maximal isometry group, then any Lorentzian manifold with constant sectional curvature $K$ carries a natural (Isom($M$), $M$)-atlas made of local isometries. \\
Simply connected space have maximal isometry group, but in general the converse is false. In particular we will focus on the case $K=-1$ and in that case it will be convenient to use non-simply connected models. 

\subsection{Basic Definitions}
By a \textit{Lorentzian metric} on a $n+1$-manifold we mean a non degenerate $2$-tensor $g$ of signature $(n,1)$. A \textit{Lorentzian manifold} is a connected manifold $M$ equipped with a Lorentzian metric $g$.\\
In a Lorentzian manifold $M$ we say that a non-zero vector $v \in TM$ is \textit{time-like} if $g(v,v)<0$, \textit{space-like} if $g(v,v)>0$ and \textit{light-like} if $g(v,v)=0$. More generally, we say that a linear subspace $V \subset T_{x}M$ is \textit{spacelike, lightlike, timelike} if the restricion of $g_x$ to $V$ is positive definite, degenerate or negative definite respectively.\\
The set of lightlike vectors, together with the null vector, disconnects $T_{x}M$ into 3 regions: two convex open cones formed by timelike vectors, one opposite to the other, and the region of spacelike vectors. 

\textcolor{red}{Inserire figura del cono di Minkowski?}

As a consequence the set of timelike in the total space $TM$ is either connected or is made by two connected components. In the latter case $M$ is said \textit{time-orientable}, and a time orientation is the choice of one of those components. Vectors is the chosen component are said     \textit{future-directed}, vectors in the other component are said \textit{past-directed}\\
A differentiable curve is said \textit{timelike, spacelike, lightlike} if its tangent vector at every point is timelike, spacelike, lightlike respectively. The curve is said to be \textit{causal} if the tangent vector is either timelike or lightlike. \\ Given a point $x$ in a time-oriented Lorentzian manifold $M$, the \textit{future} of $x$ is the set $I^+(x)$ of points which are connected to $x$ by a future-directed causal curve. The \textit{past} of $x$, $I^-(x)$, is defined in an analogus way for past-directed causal curves. 
As in the Riemannian setting, on a Lorentzian manifold $M$ there is a unique linear connection $\nabla$ which is symmetric and compatible with the Lorenzian metric $g$. We refer to it as the \textit{Levi-Civita connection} of $(M,g)$.\\ Following the analogy with Riemannian geometry the Levi-Civita connection determines the Riemann curvature tensor defined by: 
\[
    R(u,v)w=\nabla_u\nabla_v w-\nabla_v\nabla_u w-\nabla_{[u,v]}w
\]  

We then say that a Lorentzian manifold has constant sectional curvature $K$ if: 
\begin{equation}\label{sectionalcurvature}
    g(R(u,v)v,u)=K(g(u,u)g(v,v)-g(u,v)^2) 
\end{equation}
    


for every pair of vectors $u,v \in T_{x}M$ and every $x\in M$. Even tough the definition is analogus to the one given in the Riemannan realm we recall that in the Lorentzian setting the sectiona curvature can be defined only for planes in $T_{x}M$ where $g$ is non-degenerate. 

\subsection{Maximal isometry group and geodesic completeness} 

\begin{lemma}\label{isometrie} Let $M$ and $N$ be Lorentzian manifolds of constant curvature K. Then every linear isometry $L:T_{x}M\to T_yN$ extends to an isometry $f:U\to V$ where U, V are open neighbourhoods of $x,y$ respectively. Any two extensions $f:U\to V$ and $f^{\prime}:U^{\prime} \to V^{\prime} $ of $L$ coincide on $U\cap U^{\prime} $ Moreover $L$ extends to a local isometry $f:M\to N$ provided that $M$ is simply connected and $N$ is geodesically complete. 
\end{lemma}
\begin{proof}
    Cartan-Ambrose-Hicks thm.
\end{proof}
As a consequence of the aforementioned lemma we have: 
\begin{corollary}\label{122}
    Let $M$ and $N$ be simply connected, geodesically complete Lorentian manifolds of constant curvature $K$. Then any linear isometry $L:T_xM\to T_yN$ extends to a global isometry $f:M\to N.$   
\end{corollary}

In particular, there is a unique simply connected geodesically complete Lorentzian manifold of constant curvature $L$ up to isometries. For instance for $K=0$ a model is the Minkowski space $\R^{n,1}$, that is $\R^{n+1}$ provided with the standard metric: 
\[
    g=dx_1^{2}+\dots+dx_n^{2}-dx_{n+1}^2.   
\]

Another consequence of \ref{isometrie} is that, fixing a point $x_0 \in M$, the set of isometris of $M$, which we will denote by Isom($M$) can be realized as a subset of ISO($T_{x_0}M, TM$) namely the fiber bundle over $M$ whose fiber over $x\in M$ is the space of linear isometries of $T_{x_0}M$ into $T_{x}M$. \todo{Not clear what is going on here, he argues about some Lie structure.}

\begin{definition}
    A Lorentzian manifold $M$ has \textit{maximal isometry group} if the action of Isom($M$) is transitive and, for every point $x \in M$, every linear isometry $L:T_{x}M\to T_xM$ extends to an isometry of $M$. 
\end{definition}

Equivalently $M$ has maximal isometry group is the above inclusion of Isom($M$) into ISO($T_{x_0}M, TM$) is a bijection. Hence is $M$ has maximal isometry group, then the dimension of the isotropy group is maximal. 
From \ref{122}, every simply connected Lorentzian manifold $M$ has maximal isometry group if it has constant sectional curvature and is geodesically complete. The converse holds even without the simply connectedness assumption. Namely: 

\begin{lemma}\label{maximalisometry}
    If $M$ is a Lorentzian manifold with maximal isometry group, then $M$ has constant sectional curvature and is geodesically complete.
\end{lemma}

\begin{proof}
    Fix a point $x\in M$. As the identity component of $O(T_x M)\simeq O(n,1)$ acts transitebly on spacelike planes, tehre exists a constant $K$ such that \ref{sectionalcurvature} holds for every pair $(u,v)$ of vector tangent at $x$ which generate a spacelike plane. Now, for every point $x\in M$ both sides of equations \ref{sectionalcurvature} are polynomial in $u,v \in T_xM$. Since the set of pairs $(u,v)$ which generate spacelike planes is open in $T_{x}M\times T_{x}M$, the equation \ref{sectionalcurvature} must holds for every pair of vectors $u,v \in T_xM$.\\ Since Isom($M$) acts transitively on $M$, it follows that $M$ has constant sectional curvature $K$.\\
    Let us now show that the manifold is geodesically complete. Suppose $\gamma$ is a parametrized geodesic with $\gamma(0)=x$ and $\gamma^{\prime} (0)=v\in T_xM,$ which is definided for a finite maximal time $T>0.$ Let $T_0=T-\epsilon>0.$ By assumption one can find an isometry $f:M\to M$ such that $f(x)=\gamma(T_0)$ and $df_x(v)=\gamma^{\prime}(T_0).$ Then $t\to f\circ\gamma(t-T_0)$ is a parametrized geodesic which provides a continuation of $\gamma$ beyond $T$, a contradiction. 
\end{proof}

We conclude this preliminaries with a result of classification which will be useful in the following and also explains our interest in spaces with maximal isometry group.
\begin{proposition}\label{classification}
    Let $M_k$ be a simply connected Lorentzian manifold of constant sectional curvature $K$ with maximal isometry group, and let $M$ be a Lorentizian manifol of constant sectional curvature $K$. Then: 
    \begin{itemize}
        \item M is geodesically complete if and only if there is a local isometry $p:M_K\to M$ which is a universal covering.
        \item $M$ has maximal isometry group if and only if $\text(Aut)(p:M_k\to M)$ is normal in Isom($M_k$). 
    \end{itemize}
\end{proposition}

\begin{proof}
    Suppose $M$ is geodesically complete, then lifting the metric to the universal cover $\tilde{M}$ one gets a simply connected geodesically complete Lorenztian manifold of constant sectional curvature $K$ which by corollary \ref{122} is isometric to $M_K$. The covering map $p:_k\to M$ is then a local isometry by construction. The converse is straightforward. \\   
    Now let $\Gamma=\text{Aut}(p:M_K\to M),$ which is a discrete subgroup of Isom($M_K$). Thus $M$ is obtained as the quotient $M=M_{K}/\Gamma$,where $\Gamma$ acts freely and properly discontinuosly on $M_K$. The isometry group of $M$ is isomorphic to $N(\Gamma)/\Gamma$, where by $N(\Gamma)$ we denote the normalizer of $\Gamma$ in Isom($M_K$). The isomorphism is based on the observation that any isometry of $\tilde{M}$ which normalies $\Gamma$ descends to an isometry of $M$, and conversly the lifting of any isometry of $M$ must be in $N(\Gamma).$\\
    Hence the condition that $M$ has maximal isometry group is equivalent to the condition that every element $f$ of $\text{Isom}(M_K)$ descends in the quotient to an isometry of $M$. This is in turn equivalent to the condition that $f\Gamma f^{-1}=\Gamma$ for every $f\in \text{Isom}(M_K)$, which is the same as saying that $\Gamma$ is normal in $\text{Isom}(M_K)$.
\end{proof}

\begin{observation}
Since $\Gamma=\text{Aut}(p:M_K\to M)$ is discrete, being normal in $\text{Isom}(M_K)$ implies that elements of $\Gamma$ commute with the elements of the identity components of $\text{Isom}(M_K)$. This remark suggest that there are usually not many Lorenztian manifolds of constant sectional curvature with maximal isometry group.
\end{observation}

Finally, any isometry between connected open subsets of a Lorenztian manifold $M$ with maximal isometry group extends to a global isometry. In particular if $M_K$ is a Lorenziant manifold of constant sectional curvature $K$ with maximal isometry group, then any Lorentzian manifold $M$ of constant sectional curvature $K$ admits a natural $(\text{Isom}(M_K), M_K)$-structure whose charts are isometries between open subsets of $M$ and open subsets of $M_K$. We will refer to Lorentzian manifolds of constant sectional curvature $K$ with maximal isometry group as \textit{models} of constant sectional curvature $K$ (following Klein's terminology).  
\chapter{Model of Anti-de Sitter (n+1)-space}
The aim of this chapter is to construct models of Lorentzian manifolds with constant sectional curvature -1 and maximal isometry group in any dimension. We are also interested in \textit{stressing} the analogies of this manifolds with the models of hyperbolic space in the Riemannian setting. We will show that hyperbolic space is \textit{naturally} embedded in Anti-de Sitter space, and we will later develop this topic in our study of Earthquakes theory. After introducing the models, we are interested in studying the geometry of such manifolds: we will introduce a conformal boundary for Anti-de Sitter space and we will characterize geodesics and totally geodesic subspaces. We will also introduce the notion of polarity in Anti-de Sitter space (in same sense the \textit{correct} Lorentzian correspondence between points and hyperplane) and study its properties. 

\section{The quadric model}
We want to introduce the analogue of the hyperboloid model of hyperbolic space. Denote by $\R^{n,2}$ the real vector space $\R^{n+2}$ equipped with the quadratic form 
\[
    q_{n,2}(x)=x_1^{2}+\dots+x_n^{2}-x_{n+1}^{2}-x_{n+2}^2.   
\]
and by $<v,w>_{n,2}$ the associated symmetric form. Finally, let $O(n,2)$ be the group of linear transformation of $\R^{n+2}$ that preserve $q_{n,2}.$
Then we define: 
\[
    \H^{n,1}=\{x\in \R^{n,2}\;|\;q_{n,2}(x)=-1\}.
\]

It is immediate to check that $\H^{n,1},$ as the pre-image of a regular value of $q_{n,2}$, is a smooth connected submanifold of $\R^{n,2}$ of dimension $n+1$. The tangent space $T_x\H^{n+1}$, regarded as a subspace of $\R^{n+2}$ coincides with the orthogonal space $x^{\perp}=\{y\in \R^{n+2}\;|\;\langle x,y\rangle_{n,2}=0\}.$ A simple signature argument (and $q_{n,2}(x)=-1$ for every $x\in\H^{n,1}$) shows that the restriction of the symmetric form $\langle .,. \rangle_{n,2}$ to $T\H^{n+1}$ has Lorentzian signature, hence it makes $\H^{n,1}$ a Lorentzian manifold. We remark that this model is the analogue of the hyperboloid model of hyperbolic space, in fact $\H^n$ is isometrically embedded in $\H^{n,1}$ as the submanifold defined by $x_{n+2}=0$, $x_{n+1}>0$. \\
The natural action of $O(n,2)$ on $\R^{n,2}$ preserves $\H^{n,1}$, and in facts $O(n,2)$ acts by isometries on $\H^{n,1}$. We remark that $O(n,2)$ acts transitively on $\H^{n,1}$ and that the stabilizer of a point $x$ acts transitively on the space of orthonormal bases of $T_x\H^{n,1}$. Hence $\H^{n,1}$ has maximal isometry group and Isom($\H^{n,1})\simeq O(n,2)$.\\
By Lemma \ref{maximalisometry}, $\H^{n,1}$ has constant sectional curvature. Let us now check that the sectional curvature is negative (actually, $K=-1$). For this purpose, observe that the normal line in $\R^{n,2}$ to $\H^{n,1}$ at $x$ is identified with the line generated by $x$ itself. It follows that, if $v,w$ are tangent vector fields along $\H^{n,1}$ we have the orthogonal decomposition: 
\[
    (D_{v}w)(x)=(\nabla_{v}w)(x)+\langle v,w\rangle x, 
\]

where $D$ is the flat connection of $\R^{n+2}$ and $\nabla$ is the Levi-Civita connection of $\H^{n,1}$. Using the flatness of $D$ we get: 
\[
    R(u,v)w=\langle u,w\rangle v-\langle v,w\rangle u,
\]

so that 
\[
    \langle R(u,v)v,u \rangle =-(\langle u,u\rangle\langle v,v\rangle-\langle v,u\rangle^2), 
\]

and this shows that $\H^{n,1}$ has constant sectional curvature $-1$. We also remark that $\H^{n,1}$ is not simply connected, being homeomorphic to $\R^{n}\times \mathbb{S}^1$. 

\section{The ``Klein model" and it's boundary.}
Let us introduce a projective model, also knows as the ``Klein model", for Anti-de Sitter geometry. Let us define: 
\[
    \A^{n,1}=\H^{n,1}/\{\pm \text{Id}\}.
\]

Since $\{\pm \text{Id}\}$ is the center of $O(n,2)$ (hence normal), $\A^{n,1}$ (when endowed with the Lorentzian metric induced by the quotient) has maximal isometry group by Proposition \ref{classification} and is therefore a model of constant sectional curvature $-1$. It can also be shown that the center of the isometry group of the Klein model is trivial, hence $\A^{n,1}$ is the \textit{minimal} model of AdS geometry, in the sense that any other model is a covering of $\A^{n,1}$.\\
By definition $\A^{n,1}$ is naturally identified with a subspace of real projective space $\R \text{P}^{n+1}$, more explicitly with the subset of timelike directions of $\R^{n,2}$: 
\[
    \A^{n,1}=\{[x]\in \R \text{P}^{n+1}\;|\;q_{n,2}(x)<0\}.
\]
Like in hyperbolic geometry, the boundary of $\A^{n,1}$ in projective space is a quadric of signature $(n,2)$, that is the projectivization of the of lightlike vectors in $\R^{n,2}$. We denote this quadric by $\partial \A^{n,1}=\{[x]\in \R \text{P}^{n+1}\;|\;q_{n,2}(x)=0\}$.\\ We observe that isometries of $\A^{n,1}$ induce projective transformations which preserve $\partial \A^{n,1}$.  \\

\textit{The conformal Lorentzian structure of the boundary.}
We want to continue to develop analogy with hyperbolic geometry and equip $\partial \A^{n,1}$ with a conformal Lorentzian structure that extends the conformal Lorentzian structure of $\A^{n,1}$. \\
A point $\ell \in\R\text{P}^{n+1}$ is identified with $\text{Span}(x)$ for some $x\in \R^{n,2},$ and the tangent space of real projective space has the canonical identification 
\[
    T_{\ell}\R\text{P}^{n+1}\simeq \text{Hom}(\ell,\R^{n+2}/\ell). 
\]

Now if $\ell$ is timelike we can identify $\R^{n+2}/\ell$ with $\ell^{\perp}.$ For any given local section $\sigma:\A^{n,1}\to \R^{n,2}$ of the projection $\R^{n,2}\to \A^{n,1}$, we can define a Lorentzian metric on $T\A^{n,1}$ setting:

\[
    \langle \langle f,g \rangle \rangle_{\sigma}=\langle f((\sigma[x])),g(\sigma[x])\rangle_{n,2}. 
\]

for $f,g\in T_{x}\A^{n,1}\simeq \text{Hom}(\ell,\ell^\perp$). If the section $\sigma$ has image in $\H^{n,1}$, then the aforementioned metric coincides with the pull-back of the metric over $\R^{n,2}$, since the differential of $\sigma $ identifies $T_{[x]}\A^{n,1}=T_x\H^{n,1}=x^\perp$. For a general section the identification does not hold, but we can recover a conformal metric via the formula: 
\[
    \langle \langle f,g \rangle \rangle_{\lambda\sigma}=\lambda^2\langle \langle f,g \rangle \rangle_{\sigma}.
\] for any function $\lambda$.\\
We consider now the case that $\ell=\text{Span}(x)$ is lightlike, i.e. $q_{n,2}(x)=0$. In this case we can not induce any natural metric on $\R^{n,2}/\ell$. However, we let 
\[
    \mathbb{L}=\{x\in \R^{n,2}\;|\;q_{n,2}(x)=0\}
\]
be the space of lightlike vectors, then $T_x\mathbb{L}$ is precisely $\ell^{\perp}$ and contains $\ell$ itself. We have a canonical identification: $T_\ell\partial\A^{n,1}\simeq\text{Hom}(\ell,\ell^{\perp }/\ell).$ The bilinear for of $\R^{n,2},$ when restricted to $\ell^{\perp},$ induces a non degenerate bilinear form (of signature $(n-1,1)$) on $\ell^\perp/\ell.$ We will denote such a restriction as $\langle v,w\rangle_{\ell^{\perp}/\ell}.$ \\
We can now define a metric on $\partial\A^{n,1}$ for any section $\sigma:\partial\A^{n,1}\to \mathbb{L}$ of the canonical projection by the formula:
\begin{equation}\label{24}
    ((f,g))_{\sigma}=\langle f(\sigma[x]),g(\sigma[x])\rangle_{\ell^{\perp}/\ell},
\end{equation}
for all $f,g\in \text{Hom}(\ell, \ell^\perp/\ell).$ This metric can be viewed as the pull-back of the metric: 
\begin{equation}\label{25}
    ((f,g))_\sigma=\langle \sigma_*(f), \sigma_{\ast} (g)\rangle_{n,2}, 
\end{equation}
since the degenerate metric on $T_x\mathbb{L}=\ell^\perp$ is, by construction, the pull-back of the metric of $\ell^{\perp}/\ell$ by the projection along the degenerate direction $\ell$.\\ 
The relation valid for the metric on $\A^{n,1}$ also holds for the metric on $\partial\A^{n,1}$, that is: 
\begin{equation}\label{26}
    ((f,g))_{\lambda\sigma}=\lambda^2((f,g))_{\sigma}
\end{equation} 
and therefore the induced conformal class over $T\partial\A^{n,1}$ is well defined and independent of the choice of $\sigma$ and equips the tangent space of the boundary with a conformal Lorentzian metric. To conclude let us show that this conformal metric is naturally the compactification of $\A^{n,1}$. Let $\sigma$ be a section, of the projection $\pi:\R^{n,2}\to \R\text{P}^{n+1},$ defined in a neighborhood $U$ of a point in $\partial\A^{n,1}$, by construction the metric $((\cdot,\cdot))_\sigma$ over $\partial\A^{n,1}\cap U$ is the limit of the conformal metric associated to $\sigma$ defined over $\A^{n,1}\cap U$. % This means that if $(p_n,v_n)$ is a sequence in $T\A^{n,1}$ that converges to $(p_{\infty}, v_\infty)$ then 
\begin{observation}\label{222}
Let's make some observation about the light cone in the case of $\partial\A^{n,1}$. If $[y]\in\partial\A^{n,1}$ Equation \refeq{25} implies that the lightlike vectors in $T_{[y]}\partial\A^{n,1}$ are exactly the projection of vectors $x\in \R^{n,2}$ such that $\langle x,y\rangle_{n,2}=0$ and $q_{n,2}(x)=0$. These vectors are such that $\text{Span}(x,y)$ are totally degenerate planes in $\R^{n,2}$, equivalently are projective lines contained in $\partial\A^{n,1}$. The light cone in $\partial\A^{n,1}$ through $[y]$ is therefore the union of all the projective lines through $[y]$ that are contained in $\partial\A^{n,1}$.

\end{observation}


\subsection{The Poincaré model for the universal cover.}
We have already observed that $\H^{n,1}$, and its double quotient $\A^{n,1}$, are not simply connected. We want to construct a simply connected model for AdS geometry. For this purpose we introduce the universal cover of $\H^{n,1}$ and $\A^{n,1}$.\\
Let $\H^n$ be the hyperboloid model of hyperbolic space. Then: 
\begin{equation}\label{ogcover}
    \pi(y,t)=(y_1,\dots,y_n,y_{n+1}\cos t,y_{n+2}\sin t)
\end{equation}


defines a map $\pi:\H^n\times \R\to \H^{n,1}$ which is a covering with deck transformations of the form $(y,t)\mapsto (y,t+2k\pi)$ for $k\in \Z$. We denote the covering space by $\AS^{n,1}$ and we observe that it is also the universal cover of $\A^{n,1}$, where the covering map it is just the composition of $\pi$ and the double quotient.\\
Pulling back the Lorentzian metric over $\AS^{n,1}$ we get a simply connected Lorentzian manifold of constant curvature -1. The metric on $\AS^{n,1}$ is a warped product of the form: 
\begin{equation}\label{metric}
     \pi^*g_{\H^{n,1}}=g_{\H^n}-y_{n+1}^{2}dt^2.
\end{equation}
   

Moreover $\AS^{n,1}$ has maximal isometry group, hence we have obtained a simply connected model for AdS geometry. More precisely we have a central extension, that is a (non split) short exact sequence:
\[
    0\to \Z\to \textnormal{Isom}(\AS^{n,1})\to O(n,2)\to 1.
\]
It is sometimes convenient to express the metric \ref{metric} using the Poincaré model of hyperbolic space. Recall that the disk model of the hyperbolic space is the unit disk $\D^n$ endowed with the conformal metric: $\frac{4}{(1-r^2)^2}\sum dx_i^2$, where $r^2=\vert x\vert^2$. In our setting the isometry between the disk and the hyperboloid model is given by: 
\begin{equation}
    (x_1,\dots, x_n)\mapsto (y_1=\frac{2x_1}{1-r^2},\dots, y_n=\frac{2x_n}{1-r^2},y_{n+1}=\frac{1+r^2}{1-r^2}) 
\end{equation}
The Poincaré model of AdS geometry is then the cylinder $\D^n\times \R$ endowed with the metric 
\begin{equation}\label{metricdisk}
    \frac{4}{(1-r^2)^2}(dx_1^2+\dots+dx_n^2)-(\frac{1+r^2}{1-r^2})^{2}dt^2 
\end{equation}

It follows from the definition that each slice $\{t=c\}$ is a totally geodesic copy of $\H^n$. The metric defined in \refeq{metricdisk} also shows that the vector field $\partial/\partial t$ is a timelike non vanishing vector field on $\AS^{n,1}$, giving to $\AS^{n,1}$ the structure of a time-orientable manifold. Any choice of time orientation is preserved by the action of the deck transformation of the covering from the Poincaré to the Klein model, hence both $\H^{n,1}$ and $\A$ are time-orientable. \\
\begin{observation}
    Exception made for the central line passing through $x_1=\dots=x_n=t=0,$ vertical lines are not geodesic for the metric defined in \refeq{metricdisk}.  
\end{observation}

\textit{The conformal metric of the boundary.} As we did with $\A^{n,1}$ we would like to give a conformal Lorentzian structure to the boundary of $\AS^{n,1}$. The construction actually works for any covering space of $\A^{n,1}.$ The covering map (we are now considering the projective model of $\H^n$)
\[
    \pi^{\prime} ([y_1:\dots:y_{n}:y_{n+1}],t)=[y_1:\dots:y_n:y_{n+1}\cos t:y_{n+1}\sin t]
\]

extends to a map $\pi^{\prime}:\H^n\cup \partial\H^n\times \R\to \A^{n,1}\cup \partial\A^{n,1}$. In order to compute the conformal structure of the boundary, we consider the map $\tau: \H^n\times \R\to\R^{n+2}$ defined by:
\[
    \tau([y_1:\dots:y_n:y_{n+1}],t)=(y_1/y_{n+1},\dots,y_n/y_{n+1},\cos t,\sin t)
\]  
 that extends to the boundary and induces a (local) section of the projection $\R^{n,2}\to \A^{n,1}$. In fact, let $\eta$ be the generator of the group of deck transformation of the covering $\pi^{\prime}:\AS^{n,1}   \to \A^{n,1}$, then $\tau$ has the equivariance $\tau\circ\eta^i=(-1)^i\tau.$  % Using Equation \refeq{26}, the conformal Lorentzian metric on $\partial\AS^{n,1}$ induced by $\sigma$ by means of Equation \refeq{24} is the pull-back of a Lorentzian metric compatible with the natural conformal structure of the boundary $\partial\A^{2,1}$.  
\subsection{Geodesic.}

We have presented our various models as manifolds and now we would like to improve our knowledge of their geometry. As always we start by studying and characterizing geodesics.\\ 

\textit{In the quadric model} Let us start with the exponential map in the hyperboloid model. Given a point $x\in \H^{n,1}$ and a vector $v\in T_x\H^{n,1}$, we want to determine the geodesic through $x$ with speed $v$. We will distinguish several cases according to the sign of $ q_{n,2}(v)$. If $v$ is lightlike, then: 
\[
    \gamma(t)=x+tv
\] is a geodesic of $\R^{n,2}$ and is contained in $\H^{n,1}$, hence $\gamma$ is a geodesic of $\H^{n,1}$. If $v$ is etiher timelike or spacelike, we claim that the geodesic $\gamma(t)=\exp_x(tv)$ is contained in the linear plane $W=\text{Span}(x,v).$ In fact, the linear transformation $T$ that fixes pointwise $W$ and whose restriction to $W^\perp$ is $-\text{Id}_{W^\perp}$ is in $O(n,2)$. By the uniqueness of the geodesic, $T\circ\gamma=\gamma$ hence $\gamma$ is contained in $\H^{n,1}\cap W$. We can easily derive the expressions\\
\begin{equation}
    \gamma(t)=\cosh(t)x+\sinh(t)v
\end{equation}

if $q_{n,2}(v)=1$ and 
\begin{equation}\label{212}
    \gamma(t)=\cos(t)x+\sin(t)v
\end{equation}
if $q_{n,2}(v)=-1$.\\

\textit{In the Klein model.} In analogy with the Riemannian case, in the Klein model $\A^{n,1}$ geodesics are intersections of projective lines with the domain of $\A^{n,1}\subset\R \text{P}^{n+1}.$ From what we have already said: 
\begin{itemize}
    \item Timelike geodesics correspond to projective lines that are entirely contained in $\A^{n,1}$, are closed non-trivial loops and have length $\pi.$ 
    \item Spacelike geodesics correspond to lines that meet $\partial \A^{n,1}$ transversally in two points. They have infinite length.
    \item Lightlike geodesics correspond to lines tangent to $\partial\A^{n,1}$.   
\end{itemize}

In particular the light cone through a point $[x]\in \A$ coincides with the cone of lines through $[x]$ tangent to $\partial \A$.\\ For instance in the affine chart $\mathbb{A}_{n+2}\neq\{x_{n+2}=0\},$ where in coordinates $(y_1, \dots, y_{n+1})=(x_1/x_{n+2},\dots, x_{n+1}/x_{n+2})$, the intersection $\A\cap\mathbb{A}_{n+2}$ is the interior of a one sheeted hyperboloid, that is: 
\[
    \A\cap\mathbb{A}_{n+2}=\{y_1^2+\dots+y_n^2-y_{n+1}^2<1\},
\]
 while the boundary it's the one-sheeted hyperboloid itself: 
 \[
    \partial\A\cap\mathbb{A}_{n+2}=\{y_1^2+\dots+y_n^2-y_{n+1}^2<1\}.
\]
In an affine chart, timelike geodesic corresponds to affine lines which are entirely contained in the Anti de Sitter space, and which are not asymptotic to its boundary; lightlike geodesics are tangent to the one sheeted hyperboloid, or are asymptotic to it (tangent at infinity).

\begin{observation} An important observation concerns the space of timelike geodesics. Any timelike line is the projectivisation of a negative definite plane. As Isom$(\A^{n,1})\simeq \text{PO}(n,2)$ acts transitively on the space of timelike lines, and since the stabilizer of a timelike line is the group $\text{P}(O(n)\times O(2))$ which is the maximal compact subgroup of $\text{PO}(n,2)$, the space of timelike geodesics of $\A^{n,1}$ is naturally identified with the Riemannian symmetric space of PO$(n,2)$.
\end{observation}

\noindent\textit{Totally geodesic subspaces} Totally geodesic subspaces of $\A^{n,1}$ of dimension $k$ are obtained as the intersection of $\A^{n,1}$ with the projectivisation $P(W)$ of a linear subspace $W$ of $\R^{n,2}$ of dimension $k+1.$ The negative index of $W$ can be $1, 2,$ for otherwise the intersection $\A^{n,1}\cap P(W)$ would be empty. We have several cases: 
\begin{itemize}
    \item If $W$ has signature (k-1,2), then $P(W)\cap \A^{n,1}$ is isometric to $\A^{k-1,1}$. 
    \item If $W$ has signature $(k-2, 1)$ , then it is a copy of Minkowski space $\R^{k-2,1}$, hence $P(W)\cap \A^{n,1}$ is a copy of the Klein model of hyperbolic space. 
    \item If $W$ is degenerate, then $P(W)\cap \A^{n,1},$ is a lightlike subspace foliated by lightlike geodesics tangent to the same point of $\partial\A^{n,1}$. 
\end{itemize}  

A particular case of the last point is when $W$ is degenerate and $\text{dim}(W)=n+1.$ Then $P(W)\cap \A^{n,1}$ is a projective hyperplane tangent to $\partial\A$ at a point $[x]$ and $P(W)\cap \partial\A^{n,1}$ is the lightlike cone of $\partial\A^{n,1}$ through $[x]$. \todo{Check remark 2.2.2} 

\textit{In the universal cover.} In the universal cover $\AS^{n,1},$ geodesics are just the lifts of geodesics in $\A^{n,1}$ or $\H^{n,1}$. Hence every spacelike or lightlike geodesic in $\A^{n,1}$ and $\H^{n,1}$, which is topologically a line, has a countable number of lifts to $\AS^{n,1}$. Timelike geodesics in $\A^{n,1}$ and $\H^{n,1}$ are topologically circles and are in bijections with timelike geodesics in $\AS^{n,1},$ as the covering map restricted to a timelike geodesic, induces a covering map onto the circle. Using the Poincaré model we can give an explicit description of lightlike geodesic. In fact, in Lorentzian geometry not only the nature of a vector is conformally invariant but also unparametrized lightlike are a conformal invariant: 
\begin{theorem}\label{ConformalMetric} If two Lorentzian metrics $g,g^{\prime} $ on a manifold $M$ are conformal, then they have the same unparametrized lightlike geodesics.
\end{theorem}
\begin{proof}
    [GHL04, Proposition 2.131]
\end{proof}
Because of Theorem \ref{ConformalMetric} we can replace the Poincaré metric by the conformally equivalent metric given by:
\begin{equation}\label{emispherical}
    \frac{4}{(1+r^2)^2}(dx_1^2+\dots+dx_n^2)-dt^2
\end{equation} 
Now we observe that the first term in Equation \ref{emispherical} is exactly the form of the spherical metric on a hemisphere, pulled-back to the unit disk by the stereographic projection. We will call such a metric hemispherical and will denote it by $g_{\mathbb{S}^n}$. Notice that the boundary of $\partial\D$ is an equator for the hemispherical metric, and in fact it's the only equator completely contained in $(\D\cup\partial\D,g_{\mathbb{S}^n})$, a justification to the fact that we will refer to it as \textit{the} equator.\\
As a consequence, unparametrized lightlike geodesic of $\AS^{n,1}$, going through a point $(p_0,t_0)$ are characterized by the conditions that they are mapped to spherical geodesic under the vertical projection $(p,t)\to p$ and moreover: 
\[
    t-t_0=d_{\mathbb{S}^n}(p,p_0)
\] on the geodesic. In particular, these lightlike geodesic meet the boundary of $\AS^{n,1}$ at the point that satisfies the above conditions such that $p$ is on the equator of the hemisphere: as an example, if $p_0$ is the center of the hemisphere, then the points at infinity of the lightcone over $(p_0,t_0)$ are the horizontal slice $t=t_{0}+\pi/2.$ This sphere is also the boundary of a hyperplane dual to $(p_0,t_0)$, in a sense that we will explain in the following section. 

\section{Polarity in Ads}\label{polarsec}
The quadratic form $q_{n,2}$ induces a polarity on the projective space $\R\text{P}^{n+1}$, explicitly the correspondence associates to a projective subspace $P(W)$ the subspace $P(W^\perp)$. In particular, we have an induced duality between spacelike totally geodesic subspaces of $\A^{n,1}$ where the dual of a spacelike $k-$dimensional subspaces is an $n-k+1$ subspace. \\
For instance, if we consider the dual of a point $[x]\in\A^{n,1}$ it will be an $n-$dimensional spacelike hyperplane $P_{[x]}=\text{P}^\perp.$ Projectively $P_{[x]}$ is characterized as the hyperplane spanned by the intersection of $\partial\A^{n,1}$ \todo{n o n-1} with the lightcone from $[x]$. More geometrically, it can be checked that $P_{[x]}$ is the set of antipodal point to $[x]$ along timelike geodesics through $[x].$ Also, every timelike geodesic through $[x]$ meets $P_{[x]}$ orthogonally at time $\pi/2.$ Conversely, given a totally geodesic spacelike hyperplane $H$, all the timelike geodesic that meet $H$ orthogonally intersect in a single point, which is the dual point of $H$.

\subsection*{In the quadric model} We would like to \textit{lift} the idea of duality to coverings of $\A^{n,1}.$ Observe that in $\H^{n,1}$ there are two dual planes associated to any point $x$, namely the sets: 
\[
    P_x^\pm=\{\exp_x(\pm(\pi/2)v)\;|\;q_{n,2}(v)=-1,\; v\;\text{future-directed}\}.
\] 

Now the points $P_x^+$ and $P_x^-$ are antipodal and $P_{-x}^\pm=P_x^\mp$. The planes $P_x^\pm$ disconnect $\H^{n,1}$ in two regions $U_x$ and $U_{-x},$ where $U_x$ is the connected component of $x$. They can be characterized as: 
\[
    U_x=\{y\in\H^{n,1}\;|\;\langle x,y\rangle_{n,1}<0\}.
\]

Spacelike and lightlike geodesic through $x$ do not exit $U_x$, while all the timelike geodesics through $x$ meet orthogonally $P_x^\pm$ and all pass through the point $-x$. More precisely a point $y\neq x$ is connected to $x$: 
\begin{itemize}
    \item by a spacelike geodesic if and only if $\langle x,y\rangle_{n,1}<-1,$
    \item by a lightlike geodesic if and only if $\langle x,y\rangle_{n,1}=-1,$
    \item by a timelike geodesic if and only if $|\langle x,y\rangle_{n,1}|<1.$ 
\end{itemize}

An immediate consequence is that if $y$ is connected to $x$ by a spacelike geodesic, there is no geodesic joining $y$ to $-x$. Hence the exponential map of $\H^{n,1}$ is not surjective. But as any point $y\in\H^{n,1}$ can be connected through a geodesic either to $x$ or $-x$ hence the exponential over $\A^{2,1}$ is surjective.
\subsection{In the universal cover.}
Recall that $\AS^{n,1}\simeq \H^n\times\R$ with the group of deck transformations being isomorphic to $\Z$, where a generator acts by translation of $2\pi$ in the $\R$ factor. It follows that the preimage of a spacelike plane $P\subset\A^{n,1}$ is the disjoint union of spacelike planes $(P^k)_{k\in\Z},$ enumerated so that the generator $\eta$ of $\Z$ acts by sending $P^k$ to $P^{k+1}.$ \\
Each connected component of $\AS^{2,1}\setminus\bigcup_{k\in\Z} P^k$ is a fundamental domain for the action of deck transformation of the covering $\AS^{n,1}\to\A^{n,1}.$ Given a point $x$, let us apply the previous construction to the plane $P_x=P_{\pi^{\prime}(x)}$ which is the dual of the image $\pi^{\prime}(x)$ in $\A^{n,1}$, and let $V_x$ be the connected component which contains $x$. We will refer to $V_x$ as the \textit{Dirichlet domain} in $\AS^{2,1}$ centered at $x$, since the construction of $V_x$ is the analogue of a Dirichlet domain to our context. Then the restricted covering map: $\restr{\pi^{\prime}}{V_x}:V_x\to\A^{2,1}\setminus P_x$ is an isometry. Therefore lightlike and spacelike geodesic through $x$ are entirely contained in $V_x$.\\ 
\chapter{Anti-de Sitter space in dimension (2+1)}
We will now restrict our attention to dimension three Anti-de Sitter geometry, as it will be the specific model geometry of the manifolds of our next interest. In this chapter we will highlight the structure of Lie group that $\A^{2,1}$ has in this specific dimension and we will study its geometry using tools of Lie group theory. Most of the results are just a particular case of the theory developed in the previous chapter, but the Lie structure permits to give a more explicit description of the geometry of the phenomena we are interested in.\\

\section{The {PSL}$(2,\R)$ model.} 
The fundamental observation is the existence of a special model in dimension three which endows Anti-de Sitter space with a Lie group structure. To construct such a model we observe that $q=-$det is a quadratic form with signature (2,2) over the real vector space $\mathcal{M}(2,\R)$ (the signature is evident when we consider the basis consisting of elementary matrices). The associated bilinear form is expressed by the formula:
\begin{equation}\label{quadratic}
    \langle A,B\rangle=-\frac{1}{2}\text{tr}(A\cdot\text{adj}(B))
\end{equation}
for $A,B\in\mathcal{M}(2,\R)$, where adj denotes the adjugate matrix, namely: 
\[
    \text{adj}\Big(\begin{bmatrix}
        a & b \\
        c & d
    \end{bmatrix}\Big) = \begin{bmatrix}
        d & -b \\
        -c & a
    \end{bmatrix}
\]

hence, via Sylvester's theorem, there is an identification between $(\mathcal{M}(2,\R),q)$ and $(\R^{2,2}, q_{2,2}),$ unique up to composition by elements in $O(2,2)$. Under this isomorphism $\H^{2,1}$ is identified with the Lie group $\text{SL}(2,\R).$\\
Let us observe that $\text{SL}(2,\R)\times \text{SL}(2,\R)$ acts linearly on $\mathcal{M}(2,\R)$ by left and right multiplication:
$$(A,B)\cdot X=AXB^{-1}.$$

As a direct consequence of the Binet formula, the action preserves the quadratic form $q=$-det and thus induces a representation: 

\[ \rho:\text{SL}(2,\R)\times \text{SL}(2,\R)\to O(\mathcal{M}(2,\R),q). \]

Since the center of $\text{SL}(2,\R)$ is $\{\pm \text{Id}\},$ the kernel of $\rho$ is $K=\{(\text{Id},\text{Id}),(-\text{Id},-\text{Id})\},$ and by a dimensional argument (and connectedness of $\text{SL}(2,\R)$) it turns out that the image of the representation is the connected component of the identity: 
\[
    \text{Isom}_0(\H^{2,1})=\text{SO}_0(\mathcal{M}(2,\R),q)\simeq (\text{SL}(2,\R)\times \text{SL}(2,\R))/K
\]
    
Using this model, one has a natural identification between $\A^{2,1}$ and the Lie group $\PSL$, in such a way that: 
\begin{equation}\label{isomiso}
    \text{Isom}_0(\A^{2,1})\simeq \PSL\times\PSL
\end{equation}
    
acting by left and right multiplication on $\PSL$.\\ We are mostly interested in orientation-preserving and time-preserving notions that do not depend on a chosen orientation, nevertheless we will fix here an orientation and a time-orientation of $\A^{2,1}\simeq\PSL.$ As we are dealing with a Lie group it is sufficient to define on orientation of the Lie algebra, namely the tangent at the identity Id. We declare as (positive) oriented basis of $\mathfrak{sl}(2,\R):$ 
\begin{equation}\label{tangent}
V=\begin{pmatrix}
  0 & 1 \\ 1 & 0
\end{pmatrix}\;\;
W=\begin{pmatrix}
  1 & 0 \\ 0 & -1
\end{pmatrix}
\;\;
U=\begin{pmatrix}
  0 & -1 \\ 1 & 0
\end{pmatrix}
\end{equation}

The first two vectors $V,W$ are spacelike, while $U$ is timelike. $U$ is the tangent vector to the one-parameter group of elliptic isometries of $\H^2$ fixing $i\in\H^2$, parametrized by the angle of clockwise rotation; $V,W$ are vectors tangent to the one-parameter groups of loxodromic isometries fixing the geodesic with endpoints $(-1,1)$ and $(0,\infty)$ respectively. Time-orientation can also be inherited by the Lie algebra, we declare that $U$ is future-pointing timelike vector. \\

The stabilizer of the identity in Isom$_0(\A^{2,1})$ is the diagonal subgroup $\Delta<\PSL\times\PSL$. Under the obvious identification of $\PSL$ and $\Delta,$ the action of the identity stabilizer on the Lie algebra $\mathfrak{sl}(2,\R)=T_{\text{Id}}\PSL$ is the adjoint action of $\PSL$. A direct consequence of this construction is the bi-invariance of the quadratic form $q$. Indeed, denoting by $q_{\text{Id}}$ the restriction of $q$ to $T_{\text{Id}}\text{SL}(2,\R)$, a direct computation shows that $q_\text{Id}$ equals $(1/8)\kappa$ where $\kappa(X,Y)=4tr(XY)$ is the Killing form of $\mathfrak{sl}(2, \R)$.

\begin{observation}\label{311}
    The Lie algebra $\mathfrak{sl}(2,\R)$ equipped with the quadratic form $q_{\text{Id}}$ is then a copy of the 3-dimensional Minkowski space, hence the adjoint action yields a representation 
    \[
        \PSL\to O(\mathfrak{sl}(2,\R),q_{\text{Id}})
    \]

which in turn induces the well-known isomorphism: 
\[
    \text{SO}_0(2,1)\simeq\text{SO}_0(\mathfrak{sl}(2,\R), q_\text{Id})\simeq \PSL,
\] which is nothing but the restriction of the isomorphism of Equation \ref{isomiso} to the stabilizer of the identity in the left-hand side $\text{Isom}(\A^{2,1}),$ and to the diagonal subgroup $\Delta$ in the right-hand side $\PSL\times\PSL.$
\end{observation}

\section{The boundary of $\PSL$} 
From the aforementioned identification we obtain a new one between $\partial\A^{2,1}$ with the boundary of $\PSL$, inside $\text{P}(\mathcal{M}(2,\R))$ as the projectivization of the cone of rank 1 matrices. Therefore from now on we shall consider 
\[
    \partial\A^{2,1}=\{[X]\in\text{P}(\mathcal{M}(2,\R))|\;\text{rank}(X)=1\}.
\]
We endow $\overline{\A^{2,1}}\coloneqq \A^{2,1}\cup\partial\A^{2,1}$ with the topology induced by seeing both as subsets of the real projective space $\text{P}(\mathcal{M}(2,\R)).$ We want to observe that we have the following homeomorphism: 
\begin{align*}\label{delta}
    \delta:\partial\A^{2,1}&\rightarrow\R\text{P}^1\times\R\text{P}^1\\
    [X]&\mapsto (\text{Im}(X),\text{Ker}(X))
\end{align*}


where we are considering $\R\text{P}^1$ as the space of one-dimensional subspaces of $\R^2$. Since we have that $\text{Im}(AXB^{-1})=A\cdot\text{Im}(X)$ and $\text{Ker} (AXB^{-1})=B\cdot\text{Ker}(X),$ the map $\delta$ is equivariant with respect to the action of $\PSL\times\PSL,$ acting on $\partial\A^{2,1}$ as the natural extension of the group of isometries of $\A^{2,1}$ and on $\R\text{P}^1\times\R\text{P}^1$ by the obvious product action.\\
In this setting our choice of a time-orientation \textit{can be modified} according to the following Lemma:
\begin{lemma}\label{invertime}
    The inversion map $\iota[X]=[X]^{-1}$ is a time-reversing isometry of $\A^{2,1}$ which induces the homeomorphism $(x,y)\to(y,x)$ on the boundary $\partial\A^{2,1}.$
\end{lemma}
\begin{proof}
    \todo{intendono dire che $\iota[(A,B)(X)]=(\varphi(A,B))[\iota(X)]$} It follows from definition that $\iota$ is equivariant with respect to the isomorphism of $\PSL\times\PSL$ which switches left and right factors. As $\text{d}_\text{Id}\iota=-\text{Id}$ is a linear isometry, $\iota$ is an isometry, the differential being minus the identity also shows time-reversal. \\ 
    The second part of the statement can be checked via the following. For a $2\times2$ matrix the Cayley-Hamilton theorem implies the equality $(\det X)X^{-1}=(\text{tr}X)(\text{Id}-X),$ so that projectively we have $[X^{-1}]=[\text{tr}X\text{Id}-X].$ Hence $\iota$ extends to the transformation $[X]\to[\text{tr}X\text{Id}-X]$ on $\partial\A^{2,1}.$ Now if $X$ is a rank $1$ matrix, it is traceless if and only if $X^2=0$, hence $\text{Ker}(X)=\text{Im}(X).$ If $\text{tr}(X)\neq 0$, then $X$ is diagonalizable with eigenvalues $0$ and $\text{tr}(X)$. Moreover $\text{Ker}(X)$ and $\text{Im}(X)$ are the corresponding eigenspaces. It follows then that $\text{Ker}(\text{tr}X\text{Id}-X)=\text{Im}(X)$ and $\text{Im}(\text{tr}X\text{Id}-X)=\text{Ker}X$ 
\end{proof}


Consider the hyperbolic model of the upper half-plane $\H^2.$ $\R\text{P}^1$ correspond to the boundary at infinity $\partial\H^2$ and $\PSL$ is identified to $\text{Isom}_0(\H^2$). From this perspective we can consider $\partial\A^{2,1}$ as $\partial\H^2\times\partial\H^2$. We can then interpret the convergence to $\partial\A^{2,1}$ in this setting:

\begin{lemma}\label{convergenza}
    A sequence $[X_n]\in\A^{2,1}$ converges to $(x,y)\in\partial\A^{2,1}\simeq\T$ if and only if for every $p\in\H^2,$ $X_n(p)\to x$ and $X_n^{-1}(p)\to y.$ 
\end{lemma}
\begin{proof}
    $\PSL$ acts on $\H^2$ transitively via isometries, hence if the conditions holds for some $p$ then it holds for all $p\in\H^2$. Without loss of generality we can assume $p=i$ in the upper half-plane. Assuming $X_n$ converges projectively \todo{controllare se funziona} to a rank 1 matrix $X$, one checks immediately that $X(p)$ is in the projective class of $x=\text{Im}(X)$ as follows: let $X_n$ be the divergent sequence, then for any convergent sequence $v_n \in \mathbb{C}$ that does not accumulate on $\ker X$ we have that the limit of $X_n v_n$ is a non zero vector in $\text{Im}X$, so that 
    \[
        [X_n][v_n]=[X_n v_n]\rightarrow\text{Im}X\in\S.
    \]    Using this fact we have that $[A_n]\rightarrow (x,y)$ if and only if $[A_n]z\rightarrow x$ for $z\in\overline{\H^2}\setminus\{y\}$ and the convergence is uniform on compact subsets of $\overline{\H}^2\setminus\{y\}.$ \\
    The convergence of $X_n^{-1}(p)\to y$ is just a straightforward application of Lemma \ref{invertime} and an analogous argument to the one given for $X_n(p)$.
\end{proof}

In dimension three, $\partial\A^{2,1}$ is a doubly ruled quadric, which in an affine chart looks like \textcolor{red}{inserire cono}. We shall describe such rulings in a more geometric way. Given any $(x_0,y_0)\in\partial\A^{2,1}$, 
\begin{equation}\label{ruling}
        \lambda_{y_0}\coloneqq\{(x,y_0)|x\in\R\text{P}^1\}
\end{equation}

describes a projective line in $\R\text{P}^3$ which is contained in $\partial\A^{2,1}$, hence lightlike for the conformal Lorentzian structure of $\partial\A^{2,1},$ as seen in Remark \ref{222}. In fact $\lambda_{y_0}$ is the orbit of $(x_0,y_0)$ by the action of $\PSL\times\{\text{Id}\},$ or by the (now free) action of $\text{PSO}(2,\R)\times\{\text{Id}\}.$ Here $\text{PSO}(2,\R)$ corresponds to a 1-parameter elliptic subgroup in $\PSL$. In short: 
\[
    \lambda_{y_0}=\PSL\cdot(x_0,y_0)=\text{PSO}(2,\R)\cdot(x_0,y_0).
\]

We refer to $\lambda_{y_0}$ as the \textit{left ruling} through $(x_0,y_0)$, and similarly the \textit{right ruling} is 
\[
    \mu_{x_0}\coloneqq\{(x_0,y)|y\in\R\text{P}^1\}
\]

Our interest in $\partial\A^{2,1}$ is mainly justified by the following: when we will deal with earthquake theory we will often consider $\varphi:\S\to\S$ an orientation-preserving homeomorphism of the circle. The associated graph $\Lambda_\varphi$ via means of the identification given by $\delta$ is a subset of $\partial\A^{2,1}$. We observe that from equivariance of $\delta$ the following holds for every $(\alpha,\beta)\in\PSL\times\PSL$: 
\begin{equation}\label{graphequivariancy}
(\alpha,\beta)\cdot\Lambda_\varphi=\Lambda_{\beta\circ\varphi\circ\alpha^{-1}}.
\end{equation}
We remark one last time that we will consider $\partial\A^{2,1}$ as always implicitly identified with $\T$ via $\delta$.

\section{Geodesic in $\PSL$} We have already seen geodesics in the general Anti-de Sitter space, we would like to specialize here using the model of $\PSL$ and tools from general Lie Group theory. In particular we recall the following (the necessary tools in Lie groups theory are introduced in \cite{bonsanteseppi}): 
\begin{lemma}\label{331}
    Given left-invariant vector fields $V$ and $W$ on $G$, the Levi-Civita connection of a bi-invariant metric has the expression: 
    \[
        \nabla_VW=\frac{1}{2}[V,W].
    \]
    In particular, the Lie group exponential coincides with the pseudo-Riemannian exponential map.
\end{lemma}




We would like to start by considering geodesics through the identity. The Lie algebra of $\PSL$ is isometrically identified with Minkowski space, as seen in Remark \ref{311}, where under such an isometry the stabilizer of a point (namely $\PSL$ acting by means of the adjoint action) corresponds to the group of linear isometries of Minkowski space. In short, this means that we shall distinguish geodesics by their type (timelike, spacelike, lightlike) and those will be equivalent under this action. Moreover, by Lemma \ref{331} it suffices to understand the one-parameter group for the Lie group structure of $\PSL$. We immediately get the following:

\begin{itemize}
    \item Timelike geodesics are, up to conjugacy, of the form: 
    \[
    \begin{bmatrix}
    \cos(t) & -\sin(t) \\
    \sin(t) & \cos(t)
     \end{bmatrix}
\]
namely, under the identification of $\PSL$ with $\text{Isom}(\H^2)$, they are elliptic one-parameter groups fixing a point in $\H^2$. In this example, the tangent vector is the matrix
\[
    \begin{bmatrix}
    0 & -1 \\
    1 & 0
    \end{bmatrix}.
    \]
These are closed geodesics, parametrized by arclenght, of total length $\pi$. 
\item Spacelike geodesics are, again up to conjugacy, of the form: 
\[
    \begin{bmatrix}
    \cosh(t) & \sinh(t) \\
    \sinh(t) & \cosh(t) \end{bmatrix}
\]
with initial velocity: 
\[
    \begin{bmatrix}
        0 & 1 \\
        1 & 0
    \end{bmatrix}.
\]
In the hyperbolic settings, these are loxodromic one-parameter groups, fixing two points in the boundary of $\H^2$ (in this particular case, $\pm 1$). 
\item Finally, lightlike geodesics are the parabolic one-parameter groups conjugate to: 
\[
    \begin{bmatrix}
        1 & t \\
        0 & 1 
    \end{bmatrix},
\]
whose initial vector has indeed zero length.
\end{itemize}

\subsection{Timelike geodesics}
When looking for a complete description of timelike geodesics it suffices to let (the identity component of) the isometry group of $\A^{2,1}$ (namely $\PSL\times\PSL)$ act on $\PSL$ via left and right multiplication. In particular, we can describe the whole space of timelike geodesics of $\A^{2,1}$ as follows: 

\begin{proposition}\label{352}
    There is a homeomorphism between the space of (unparametrized) timelike geodesics of $\A^{2,1}$ and $\H^2\times\H^2$. The homomorphism is equivariant for the action of $\text{Isom}_{0} (\A^{2,1})\simeq \PSL\times\PSL$.
\end{proposition}
\begin{proof}
    The homeomorphism is defined as follows. Given a point $(p,q)\in \H^2\times\H^2,$ we associate to it the subset: 
    \[
        L_{p,q}=\{X\in\PSL\;|\;X\cdot q=p\}
    \]
    By the previous discussion, timelike geodesics through the identity are precisely of the form $L_{p,p}$ for some $p\in\H^2.$ The map $(p,q)\mapsto L_{p,q}$ is equivariant for the natural actions of $\PSL\times \PSL$, namely $(A,B)\cdot L_{p,q}=L_{A\cdot p,B\cdot q},$ which also implies that $L_{p,q}$ is an unparameterized timelike geodesic and that all the unparameterized timelike geodesics are of this form, namely the map we defined is surjective. It remains to show injectivity; if $L_{p,q}=L_{p^{\prime} ,q^{\prime} }$ for $(p,q)\neq(p^{\prime},q^{\prime} )$ then in particular there exists an isometry $X_1$ of $\H^2$ sending $p$ to $q$ and $p^{\prime} $ to $q^{\prime}$, but such an isometry is necessarily unique. Suppose the existence of an $X_2\neq X_1$ isometry of $\H^2$ with the same property, then $X_2^{-1}\circ X_1$ fixes $p,p^{\prime}$, an absurd since the identity is the only isometry of $\H^2$ fixing two different points.
\end{proof}

\subsection{Spacelike geodesics}
Let us consider $\ell$ an oriented geodesics of the hyperbolic plane $\H^2$. From what we have already discussed the one-parameter group of loxodromic transformation fixing $\ell$ as an \textit{oriented} geodesic constitutes a spacelike geodesic through the origin. By an argument analogous to the one given in Proposition \ref{352}, relying on the equivariance of the construction by the actions of $\PSL\times\PSL$, one proves that every spacelike geodesic is of the form: 
\[
    L_{\ell,\jmath}=\{X\in\PSL\;|\;X\cdot\jmath=\ell\;\text{as oriented geodesics}\},
\]
where $\ell$ and $\jmath$ denote oriented geodesics of $\H^2$. We want to emphasize that every (unparameterized, unoriented) spacelike geodesic can be expressed in the above form in two ways, as we could change the orientation of both $\ell$ and $\jmath$. Rephrasing we can state: 
\begin{proposition}\label{353}
    There is a homeomorphism between the space of (unparametrized) oriented spacelike geodesics of $\A^{2,1}$ and the product of two copies of $\partial\H^2\times\partial\H^2\setminus\Delta,$ the space of oriented geodesics of $\H^2$. The homeomorphism is equivariant for the action of $\text{Isom}_0(\A^{2,1})\simeq\PSL\times\PSL$.
\end{proposition}

However we will not be really interested in oriented geodesics, hence we will have identification $L_{\ell,\jmath}=L_{\ell^{\prime},\jmath^{\prime}}$ where with $\ell^{\prime} $ we denote $\ell$ endowed with the opposite orientation.\\
Given a spacelike geodesic, there is a natural notion of \textit{dual} spacelike geodesic, which we define using the projectivity duality between points and planes discussed in section \ref{polarsec}: 
\begin{definition}
    Given a spacelike geodesic $L_{\ell,\jmath}$ in $\A^{2,1}$, the \textit{dual spacelike geodesic} is the intersection of all the spacelike planes dual to points of $L_{\ell,\jmath}.$  
\end{definition}

The construction of the dual geodesic is an involution, let us see it with an explicit example. Consider the geodesic $L_{\ell,\ell}$ through the origin, which consists of the one-parameter loxodromic group of $\PSL$ translating along $\ell$, it can be checked that the dual geodesic consists of all elliptic order-two elements (hence contained in $P_{\text{Id}}$, as we could expect from the definition) whose fixed point lies in $\ell$. To see this we can suppose (as a consequence of the transitivity of the action of $\PSL$ on geodesic of $\H^2$ ) that $\ell$ is the imaginary axis. In this case $L_{\ell,\ell}=\{M_k|k\in\R\}$ where \[
    M_k=\begin{bmatrix}
        e^{-k/2} & 0 \\
        0 & e^{k/2} \\
    \end{bmatrix}.
    \]
It follows from \refeq{geoplanes} that $X\in P_{[M_k]}$ if and only if $\langle X, M_k\rangle=0$, which given a generic matrix
\[
    X=\begin{bmatrix}
        a & b \\

        c & d
        \end{bmatrix}
\]

can be restated as the condition \begin{equation}\label{ellittichine} 
    ae^{\frac{k}{2}}+de^{\frac{-k}{2}}=0.
\end{equation}
As we are asking that Equation \refeq{ellittichine} holds for every $k$ it must be $a=d=0$, hence $X$ (seen as a Möbius transformation) is of the form $X(z)=-\frac{b^2}{z}$ an elliptic transformation of order two with fixed point $ib\in \H^2$.\\
In other words, the dual spacelike geodesic of $L_{\ell,\ell}$ is $L_{\ell,\ell^{\prime}}.$\\
We can explicitly describe the points at infinity in $\partial\A^{2,1}$ of these geodesics. Using Lemma \ref{convergenza}, if $x$ and $y$ are the endpoints at infinity of $\ell$ in $\partial\H^2,$ then any sequence diverging towards an end of $L_{\ell,\ell^{\prime}}\subset \PSL$ maps an interior point towards $x$, and the sequence of inverses towards $y$ (up to switching the two points). In other words, under the identification given by $\delta$ between $\partial\A^{2,1}\simeq\T$, the endpoints of $L_{\ell,\ell}$ are $(x,y)$ and $(y,x)$. A similar argument applied to the geodesic $L_{\ell,\ell^{\prime} },$ which consists of order-two elliptic isometries with fixed point in $\ell$ shows that its endpoints are $(x,x)$ and $(y,y).$\\
Recall the description of the left and right rulings of $\partial\A^{2,1}$ given in \refeq{ruling}, we can conclude that the endpoints of a spacelike geodesic and its dual are mutually connected by lightlike segments in $\partial\A^{2,1}$. \textcolor{red}{Qua andrebbero inserite figure per referenza.}\\
The naturality of the above construction with respect of the action of $\PSL\times\PSL$ let us state: 
\begin{proposition}
    Given a spacelike geodesic $L_{\ell,\jmath}$ of $\A^{2,1}$, its endpoints in $\partial \A^{2.1}$ are $(x_1,y_2)$ and $(y_1,x_2)$, where $x_1$ and $y_1$ are the final and initial endpoints of $\ell$ in $\partial\H^2$, and $x_2$ and $y_2$ are the final and initial endpoints of $\jmath$. The dual geodesic is $L_{\ell,\jmath^{\prime}}$ and has endpoints $(x_1,x_2)$ and $y_1,y_2.$
\end{proposition}


\section{Spacelike planes}
Now we want to study totally geodesic spacelike planes in $\A^{2,1}.$ They are all obtained as the intersection of $\A^{2,1}$ with a projective subspace in the projective space $\text{P}\mathcal{M}(2,\R).$ Hence they are all of the the following form: 
\begin{equation}\label{geoplanes}
    P_{[A]}=\{[X]\in\PSL\;|\;\langle X,A\rangle=0\}
\end{equation}
for some non-zero matrix $A$. The notation is justified by the observation that the plane defined by $P_A$ depends only on the projective class of $A$. The totally geodesic plane is spacelike if and only if $q(A)=-\det A$ is negative. We will call such a plane \textit{dual plane} of $A$, in particular the dual plane $P_\gamma$  of an element $\gamma\in\PSL$ is a spacelike totally geodesic plane. \\
\begin{example}\label{spacelikeisometry}
Before the general treatment we want to focus on a concrete example. Consider $\gamma=\text{Id}\in\PSL.$ Now because of Equation \refeq{quadratic}, $P_{\text{Id}}$ is the subset of $\PSL$ consisting of projective class of unit matrices $X$ with $\text{tr}(X)=0$. By the Cayley-Hamilton theorem, $X^2=-\text{Id},$ hence the elements of $P_{\text{Id}}$ are order-two isometries of $\H^2,$ that is, elliptic elements with rotation angle $\pi$.  We can also observe that $P_{\text{Id}}$ is invariant under the actions of $\PSL$ by conjugation, which corresponds to the diagonal in the isometry group $\PSL\times\PSL$ of $\A^{2,1}$. Using Lemma \ref{convergenza}, we can see that the boundary of $P_{\text{Id}}$ in $\partial\A^{2,1}\simeq \R\text{P}^1\times\R\text{P}^1$ is the diagonal; more precisely: 
\begin{equation}\label{graphino}
    \partial P_{\text{Id}}=\text{graph(Id)}\subset \T.
\end{equation}

Now consider a point $z\in\H^2$, and let us denote by $\mathcal{R}_z$ the order-two elliptic isometry with fixed point $z$. We claim that the map 
\[
    \iota:\H^2\to P_{\text{Id}},\;\;\;\ \iota(z)=\mathcal{R}_z
\]

is an isometry with respect to the hyperbolic metric of $\H^2$ and the induced metric on $P_{\text{Id}}.$ First, the inverse of $\iota$ is simply the fixed-point map $\text{Fix}:P_{\text{Id}}\to\H^2$ sending an elliptic isometry to its fixed point, which also shows that $\iota$ is equivariant with respect to the action of $\PSL$ on $\H^2$ by homographies and on $P_\text{Id}$ by conjugation, since $\text{Fix}(\alpha\gamma\alpha^{-1})=\alpha\text{Fix}(\gamma).$ That is, the following holds: 
\begin{equation}
    \iota(\alpha\cdot p)=\alpha\circ\iota(p)\circ\alpha^{-1}.
\end{equation} 
A direct consequence is that $\iota$ is isometric, since the pull-back of the metric of $P_\text{Id}$ is necessarily $\PSL$-invariant and has constant curvature -1, hence it coincides with the standard hyperbolic metric on the upper half-plane.
\end{example}
This simple example is actually the key case to understand general spacelike totally geodesic planes as every spacelike totally geodesic plane is of the form $P_\gamma$ for some $\gamma\in\PSL.$ To see this, observe that the action of the isometry group of $\A^{2,1}$ on spacelike totally geodetic planes is transitive, and that $P_\gamma=(\gamma,\text{Id})\cdot P_{\text{Id}}$ as the isometry $(\gamma,\text{Id})$ maps $\text{Id}\to\gamma$, and therefore the dual plane of $\text{Id}$ to the dual plane of $\gamma.$ In view of the observations given in \refeq{graphequivariancy} and \refeq{graphino}, we can conclude the following: 
\begin{lemma}\label{32}
    Every spacelike totally geodesic plane of $\A^{2,1}$ is of the form $P_\gamma$ for some orientation-preserving isometry $\gamma\in\PSL,$ and 
    \[
        \partial P_\gamma=\text{graph}(\gamma^{-1})\subset\T.
    \]
\end{lemma}
\section{Timelike planes}
Let us now consider a matrix $A\in\mathcal{M}(2,\R)$ such that $\det(A)=-1.$ The corresponding plane $P_A$ defined by equation \refeq{geoplanes} is a timelike totally geodesic plane. Associated with $[A]$ is an orientation-reversing isometry $\eta$ of $\H^2.$ Moreover, the action of $A$ by homography on $\mathbb{C}\text{P}^1$ preserves $\S$ and switches the two connected components of the complement. The matrix $A$ thus induces an orientation-reversing isometry, up to identifying the two components via $z\to\overline{z}$. We will thus denote $P_{[A]}$ by $P_\eta$. \\
The totally geodesic timelike plane $P_\eta$ can now be parametrized as follows. We have a map: 
\begin{equation}\label{refspa}
    \mathcal{I}\to\mathcal{I}\eta
\end{equation}
from the spaces of reflections $\mathcal{I}$ along geodesic of $\H^2$, with values in $\PSL\simeq\A^{2,1}$. Its image it precisely $P_\eta.$ As seen in the proof of \ref{invertime} we have that an $X$ with determinant $-1$ is an inversion if and only if traceless. As $\det(A)=-1$ we have $\text{adj}(A)=-A^{-1}$, therefore $\langle XA,A\rangle=0$ if and only if $\text{tr}(X)=0$, that is $X$ is traceless hence an involution. This shows that the image of the map defined in Equation \refeq{refspa} is the entire plane $P_\eta$.\\  
In similar fashion to the spacelike case, using the transitivity of the group of isometries on timelike planes, every timelike planes is of the form above. We can show the following: 
\begin{lemma}\label{33}
    Every timelike totally geodesic plane of $\A^{2,1}$ is of the form $P_\eta$ for some orientation-reversing isometry $\eta\in\PSL,$ and 
    \[
        \partial P_\eta=\text{graph}(\eta^{-1})\subset\T.
    \]
\end{lemma}
\begin{proof}
    We want to use Lemma \ref{convergenza} and the parametrization given in Equation \refeq{refspa}. Suppose we have a sequence $\mathcal{I}_n$ such that $\mathcal{I}_n\eta(z)\to x\in\partial\H^2$, for any $z\in\H^2$. Then, using that $\mathcal{I}_n$ is an involution and the continuity of the action of $\eta$ on $\overline{\H}^2$, $(\mathcal{I}_n\eta)^{-1}(z)=\eta^{-1}\mathcal{I}_n^{-1}(z)=\eta^{-1}\mathcal{I}_n(z)=\eta^{-1}(x).$ 
\end{proof}
\begin{observation}
    Since reflections of $\H^2$ are uniquely determined by (unoriented) geodesics, we can consider the map defined in \refeq{refspa} as a map from the space $\mathcal{G}(\H^2)$ of unoriented geodesics of $\H^2$ to $\PSL.$ This map is isometric with respect to the natural metric on $\mathcal{G}(\H^2)$ which makes it identified with the two-dimensional space $\A^{1,1}$ as treated more in detail by El Emam e Seppi in \cite{emam2022gauss} (Example $6.1$).
    \end{observation}
\section{Lightlike planes}
We are only left with case of lightlike totally geodesic planes. Those are the form $P_{[A]}$ for a nonzero matrix $A$ with $\det(A)=0$. Before giving an explicit description of such planes, we want to observe that their boundary will not be a graph in $\T$, unlike the case of spacelike and timelike planes. 

\begin{lemma}\label{35}
    Every lightlike totally geodesic plane of $\A^{2,1}$ is of the form $P_{[A]}$ for some rank one matrix $A$, and: 
    \[
        \partial P_{[A]}=(\text{Im}(A)\times\S)\cup(\S\times\text{Ker}(A)).
    \]
\end{lemma}
\begin{proof}
    The point in $\partial P_{[A]}$ are projective class of rank one matrices satisfying $\langle X,A\rangle=0$, that is, such that $\text{tr}(X\text{adj}(A))=0.$ Given that $X\text{adj}(A)$ has vanishing determinant, by the Cayley-Hamilton theorem $X\text{adj}(A)$ is traceless if and only if it is nilpotent, that is, if and only if $X\text{adj}(A)X\text{adj}(A)=0.$ Now, given that image and kernel of both $X$ and $\text{adj}(A)$ have dimension 1, this happens if and only if: 
    \begin{equation}
        \text{Im}(X)=\text{Ker}(\text{adj}(A))\;\text{or}\;\text{Im}(\text{adj}(A))=\text{Ker}(A).
    \end{equation}
    Now, since $\det(A)=0$ implies $\text{adj}(A)A=A\text{adj}(A)=0$, the relations $\text{Ker}(\text{adj}(A))=\text{Im}(A)$ and $\text{Im}(\text{adj}(A))=\text{Ker}(A)$ hold. Hence $X\in P_{[A]}$ if and only if $\text{Im}(X)=\text{Im}(A)$ or $\text{Ker}(X)=\text{Ker}(A),$ which concludes because of the definition of the homeomorphism $\delta$.  
\end{proof}

Geometrically $\partial P_{[A]}$ is the union of two circles in $\T,$ one horizontal and one vertical, which intersect exactly at the point in $\T$ corresponding to $[A]\in \partial\A^{2,1}$ via $\delta$.\\


\chapter{Mess Work}

In his 1990 paper "Lorentz Spacetimes of Constant Curvature" \cite{Mess}, Geoffrey Mess offered a completely new approach to the study of spacetimes in 2+1-dimension by employing tools and techniques from low-dimensional geometry and topology. The aim of this chapter is to give a brief introduction to Mess ideas, with a special attention to AdS geometry. Our treatment will follow the setting introduced by Bonsante and Seppi in \cite{bonsanteseppi}. We want to introduce general Lorentzian sets such as \textit{achronal subset, invisible domain, domain of dependence} and how they relate to graph of circle homeomorphism (and its convex hull) in our specialized Anti-de Sitter setting. We will show that those graphs are \textit{proper achronal sets} in the projective model and always lift to achronal sets in the Poincaré model.

\section{Causality and Convexity properties}
We begin by giving some definitions:
\begin{definition}
    A subset $X$ of $\AS^{2,1}\cup \partial\AS^{2,1}$ is achronal (respectively acausal) if no pair of points in $X$ is connected by timelike (resp. causal) lines in $\AS^{2,1}$.
\end{definition}
Since acausality and achronality are conformally invariant notions, it will be often convenient to consider the metric $g_{\mathbb{S}^2}-dt^2$ on $\D\times\R$ introduced in \refeq{emispherical} which is conformal to the Poincaré model. We give now a first useful characterization of achronal and acausal sets:
\begin{lemma}
    A subset $X$ of $\AS^{2,1}\cup \partial\AS^{2,1}$ is achronal (respectively acausal) if and only if it is the graph of a function $f:D\to\R$ which is 1-Lipschitz, (resp. strictly 1-Lipschitz) with respect to the distance induced by the hemispherical metric $g_{\mathbb{S}^2}$.\\
    Where we have denoted $D=\pi_{\D}(X).$ 
\end{lemma}
\begin{proof}
     Let's assume $X$ is achronal. Now, since vertical lines in the Poincaré model are of timelike type, the restrictions of the projections $\pi_\D:\D\times\R\to\D$ to $X$ are injective. But then, $X$ can be interpreted as the graph of a function $f:D\to\R$. By imposing that $(x,f(x))$ and $(y,f(y))$ are not connected by a timelike curve we deduce that: 
     \begin{equation}\label{soloqua}
        |f(x)-f(y)|\leq d_{\mathbb{S}^2}(x,y)
     \end{equation}
     where $d_{\mathbb{S}^2}$ is the distance induced by the hemispherical metric. By the same reasoning we show that a 1-lipschitz graph over $\D$ is achronal. Moreover two points $(x,t)$ and $(y,s)$ are on the same lightlike geodesic if and only if $d_{\mathbb{S}^2}(x,y)=|t-s|$. Hence $X$ is acausal if and only if the inequality \refeq{soloqua} is strict. \
\end{proof}

Now a 1-Lipschitz function on a region $D\subset \D$ extends uniquely to the boundary of $D$. As a simple consequence of the previous lemma, we thus have: 

\begin{lemma}\label{achronalgraph}
    An achronal subset $X$ in $\AS^{2,1}$ is properly embedded if and only if it is a global graph over $\D$, and in this case it extends uniquely to the global graph of a 1-Lipschitz function over $\D\cup\partial\D$.
\end{lemma}

Because of Lemma \ref{achronalgraph} we will often refer to an \textit{achronal surface} as an achronal subset $X\subset\AS^{2,1}$ which is the graph of a 1-Lipschitz function defined in a domain in $\D.$ Before moving over to the study of properties we shall remark how achronality and acausality are global conditions. 

\begin{definition} Given a surface $S$ and a Lorentzian manifold $(M,g),$ a $\mathcal{C}^1$ immersion $\sigma:S\to M$ is \textit{spacelike} if the pull-back metric $\sigma^*g$ is Riemannian. If $\sigma$ is an embedding, we refer to its image as a \textit{spacelike surface}.
\end{definition}

A spacelike surface $S$ is locally acausal, but there are examples of spacelike surfaces which are not achronal (hence a fortiori not acausal), a fact that highlights the global character of the achronality condition. On the other hand the following is true: 

\begin{lemma}
    Any properly embedded spacelike surface in $\AS^{2,1}$ is acausal. 
\end{lemma}
\begin{proof}
    By Lemma \ref{achronalgraph}, any properly embedded spacelike surfaces $S$ in $\AS^{2,1}$ disconnects the spaces in two regions $U,V$ whose common boundary is $S$, and we can assume that the outward pointing normal from $U$ (resp. $V$) is past-directed (resp. future directed). It then turns out that any future oriented causal path that meets $S$ passes from $V$ towards $U$. This implies that any causal path meets $S$ at most once.
\end{proof}

We have remarked in Theorem \ref{ConformalMetric} that unparametrized lightlike geodesic only depend on the conformal class of the Lorentzian metric, hence we will just refer to lightlike in $\AS^{2,1},$ even when we are considering it endowed with the hemispherical metric of \refeq{emispherical}. 


\begin{lemma}\label{containedgeo}
    Let $S$ be a properly embedded achronal surface of $\AS^{2,1}\cup\partial\AS^{2,1}$ and assume that a lightlike geodesics segment $\gamma$ joins two points of $S$. Then $\gamma$ is entirely contained in $S.$
\end{lemma}
\begin{proof}
    We want to exploit Lemma \ref{achronalgraph}, let $f^S:\overline{\D}\to\R$ be the function that defines $S$, which is $1$-Lipschitz with respect to the hemispherical metric. Now if our segment $\gamma$ joins $(x,f^S(x))$ to $(y,f^S(y))$, then (up to switching the role of $x$ and $y$) it holds: $f^S(y)=f^S(x)+d_{\mathbb{S}^2}(x,y).$ More is true; if $\gamma$ consists of points of the form $(z,f^S(x)+d_{\mathbb{S}^2}(x,z))$, for $z$ that lie on the $g_{\mathbb{S}^2}$ geodesic connecting $x$ to $y$. For such a point $z$ on the geodesic, because of achronality of $S$, it holds: 
        \[
            f^S(z)-f^S(x)\leq d_{\mathbb{S}^2}(x,z)\;\;\text{and}\;\;f^S(y)-f^S(z)\leq d_{\mathbb{S}^2}(z,y)=d_{\mathbb{S}^2}(x,y)-d_{\mathbb{S}^2}(x,z).
        \]
        Second inequality implies that $f^S(z)\geq f^S(x)+d_{\mathbb{S}^2}(x,z),$ it follows that $f^S(z)=f(x)+d_{\mathbb{S}^2}(x,z)$ and hence $\gamma$ is entirely contained in $S.$
\end{proof}

\subsection{Invisible domain}
Invisible domain were not treated in the original Mess work but were later introduced by Barbot in the literature in \cite{barbot2008causal}. We give general definition and properties for a generic $X$ subset of $\AS^{2,1}\cup\partial  \AS^{2,1}$ and then we will focus on $X$ entirely contained in the boundary. 

\begin{definition}
    Given an achronal domain $X$ in $\AS^{2,1}\cup\partial\AS^{2,1}$, the \textit{invisible domain} of $X$ is the subset, that we will denote by $\Omega(X),$ of $\AS^{2,1}$ defined as the set of points which are not connected to $X$ by no causal path.   
\end{definition}

We recall that any 1-Lipschitz function on a subset of a metric space admits a $1-$Lipschitz extension everywhere (Mc Shane's theorem \cite{mcshane1934extension}). In our setting this allows us show that any achronal set $X$ is a subset of a properly embedded achronal surface.\\
We consider two particular extensions $f_\pm^X:\D\cup\partial\D,$ which we will refer to as \textit{extremal} extensions:
\begin{align*}
    f_-^X(y)&=\sup\{f^X(x)-d_{\mathbb{S}^2}(x,y)\;|\;x\in\pi_\D(X)\}    \\
    f_+^X(y)&=\inf\{f^X(x)+d_{\mathbb{S}^2}(x,y)\;|\;x\in\pi_\D(X)\}.
\end{align*}
    
\begin{lemma}\label{422}
    Let X be any closed achronal subset $X$ of $\AS^{2,1}\cup\partial\AS^{2,1}$, and let $S_\pm(X)$ be the graphs of the extremal extensions $f_\pm^X$. 
    \begin{itemize}
        \item The properly embedded surfaces $S_-(X)$ and $S_+(X)$ are achronal with $S_-(X)\subset \overline{I^-(S_+(X))}$, and $\Omega(X)=I^+(S_-(X))\cap I^-(S_+(X)).$
        \item Every achronal subset containing $X$ is contained in $S_-(X)\cup\Omega(X)\cup S_+(X).$
        \item Every point of $S_\pm(X)$ is connected to $X$ by at least one lightlike geodesic segment, which is contained in $S_\pm(X).$ Finally $S_+(X)\cap S_-(X)$ is the union of $X$ and all lightlike geodesic segment joining points of $X$.
    \end{itemize}
\end{lemma}
\begin{proof}
    We begin by showing that $S_-(X)\subset \overline{\text{I}^-(S_+(X))}$. Given a point $(y,t),\;t\leq f_+^X(y)$ if and only if $t\leq f^X(x)+d_{\mathbb{S}^2}(x,y)$ for every $x\in\pi_{\D}(X),$ that is, if and only if $(y,t)$ lies outside $\text{I}^+(X)$. Similary $(y,t)$ lies outside $\text{I}^-(X)$ if and only if $t\geq f_-^X(y).$ By achronality, $S_+(X)$ does not meet the past of $X$, so we deduce that $f_+^X(y)\geq f_-^X(y)$ for all $y\in\overline{\D}$. Hence $S_-(X)\subset \overline{\text{I}^-(S_+(X))}$.\\
    As a similar observation, given a point $(y,t)$ we have that $\{(y,t)\}\cup X$ is achronal if and only if $f_-^X(t)\leq t\leq f_+^X(y)$. Moreover $(y,t)$ is connected to $X$ by no causal curve if and only if $f_-^X(y)<t<f_+^X(y)$. THis shows that 
    \[
        \Omega(X)=\{(y,t)\;|\;f_-^X(y)<t<f_+^X(y)\}.
    \]
    and also the second item, by applying the previous observation to any point of an achronal set containin $X$ which is not in $X$ itself. \\
    To prove the third item, fix a point $(y,t)\in S_+(X)$. As we are assuming that $X$ is closed in $\AS^{2,1}\cup\partial\AS^{2,1}$, the fact that $f^X$ is 1-Lipschitz implies that $\pi_{\D}(X)$ is closed in $\D\cup\partial\D$, so it is also compact. In particular, there exists $x\in\partial\D$ such that $t=f_+^X(y)?f^X(x)+d_{\mathbb{S}^2}(x,y)$. Thus $(y,t)$ is connected to $(x,f^X(x))$ by a lightlike geodesic segment. By Lemma \ref{containedgeo} this geodesic is entirely contained in $S_+(X)$. Clearly the same proof works for $S_.(X)$. \\
    We are left with the computation of $S_-(X)\cap S_+(X)$. For this purpose, notice that if two points of $X$ are connected by a lightlike goedesic segment, again via Lemma \ref{containedgeo}, we deduce that $\gamma\subset S_-(X)\cap S_+(X).$ Converely let $(y,t)\in S_-(X)\cap S_+(X)$ so that $f_-^X(y)=f_+^X(y).$ There exist $x$ and $x^{\prime}$ in $\pi_{\D}(X)$ such that 
    \[
        f_+^X(y)=f^{X}(x)+d_{\mathbb{S}^2}(x,y) \;\;\text{and}\;\;f_-^X(y)=f^X(x^{\prime})-d_{\mathbb{S}^2}(x^{\prime},y). 
    \]
    Using the equality $f_-^X(y)=f_+^X(y),$ the triangle inequality and the fact that $f^X$ is 1-Lipschitz we deduce that 
    \begin{equation}\label{42}
        f^{X}(x)-f^X(x^{\prime})=d_{\mathbb{S}^2}(x,x^{\prime})=d_{\mathbb{S}^2}(x,y)+d_{\mathbb{S}^2}(y,x^{\prime}). 
       \end{equation}
       Hence the points $(x,f^X(x))$ and $(x^{\prime},f^X(x^{\prime} ))$ are joined by a lightlike segment. If $x,x^{\prime}$ are not antipodal points on $\partial\D$, there is a unique hemispherical geodesic $\eta$ in $\overline{\D}$ joining $x$ to $x^{\prime}$, which must pass through $y$ by \refeq{42}, and which we may assume parametrized by arcleght. In this case the geodesic segment joining $(x,f^X(x))$ to $(x^{\prime},f^X(x^{\prime} ))$ take the form $t\mapsto(\eta(t),f^X(x^{\prime})+t)$, so it passes through $(y,f_+^X(y))=(y,f_-^X(y)).$\\
       If $x$ and $x^{\prime}$ are antipodal, then there are infinitely many geodesic joining $x$ to $x^{\prime}$, and we can pick one going through $y$. Then the same argument as above applies. 
\end{proof}

\begin{observation}\label{423}
    Given a point $(y,t)$ the set of points $(x,s)$ satisfying $\lvert s-t \rvert<d_{\mathbb{S}^2}(x,y)$ coincides with the region of $\AS^{2,1}$ which is connected to $(y,t)$ by a spacelike geodesic for the Anti-de Sitter metric. It coincides also with the region of points connected to $(y,t)$ by a spacelike geodesic for the conformal emispherical metric. Now, since $f_-^X(y)\leq t\leq f_+^X(y)$ is equivalent to the condition that $\lvert s-t \rvert<d_{\mathbb{S}^2}(x,y)$ for all $(x,t)\in X,$ the region 

    \[
        S_+(X)\cup\Omega(X)\cup S_-(X)=\{(y,t)\;|\; f_-^X(y)\leq t\leq f_+^X(y)\}
    \]
    consist of all the points that are connected to any point of $X$ by spacelike or lightlike geodesics. Moreover $\Omega(X)$ consist of points connected to any point of $X$ by a spacelike geodesic. We observe that $\Omega(X)$ could be empty, for instance when $X$ is a global graph then $S_-(X)=S_+(X)=X$ and $\Omega(X)$ is empty.
\end{observation}

\begin{observation}
    Since any point of $S_\pm(X)$ is connected to $X$ by a lightlike geodesic, it follows from Lemma \ref{containedgeo} that the intersection of any properly embedded achronal surface containing $X$ with $S_\pm(X)$ is a union of lightlike geodesic segments with an endpoint in $X$. In particular any properly embedded acausal surface containing $X$ is contained in $\Omega(X).$
\end{observation}

We will need the following technical definition: 

\begin{definition}
    Given a function $f:\overline{\D}\to\R$, we define its oscillation as: 
    \[
        \text{osc}(f)\coloneqq \text{max}_{y\in\overline{\D}}f(y)-\min_{y\in\overline{\D}}f(y).
    \]
    We observe that such a quantity is \textit{not} invariant under the action of the isometry group of $\AS^{2,1}.$
\end{definition}

With this definition we can state the following lemma that we will soon use: 

\begin{lemma}\label{417}
    Let $S$ be a properly embedded achronal surface, defined as the graph of $f^S:\overline{\D}\to\R$. Then $\text{osc}(f^S)\leq\pi$ moreover $\text{osc}f^S=\pi$ if and only if $S$ is a lightlike plane. 
\end{lemma}
\begin{proof}
    We observe that $f^S$ is $1-$Lipschitz for the hemispherical metric, and the diameter of $\D$ for $g_{\mathbb{S}^2}$ is $\pi$, hence we easily obtain that $\text{osc}(f^S)$ is bounded by $\pi$. Moreover if the value is attained it follows that there are two antipodal points $y,y^{\prime} \in \partial\D$ such that: $f^S(y^{\prime})=f^S(y)+\pi$. We have already notice that the lightlike plane with past and future points $(y,f^S(y))$ and $(y^{\prime},f^S(y)+\pi)$ is: 
    \[
        P=\{(x,t)\;|\;t=f^S(x)+d_{\mathbb{S}^2}(x,y)\}
    \]
    and it is foliated by lightlike geodesic joining $(y,f^S(y))$ to $(y^{\prime},f^S(y)+\pi)$. By Lemma \ref{containedgeo}, $P$ is included in $S$. Since both are global graph over $\overline{\D}$ we have $S=P$.
\end{proof}

\subsection{Achronal meridian in $\partial\AS^{2,1}$}
We will be interested in the study of invisibile domains of achronal meridians $\Lambda$ in the boundary of $\AS^{2,1},$ that are graphs of 1-Lipschitz functions $f:\partial\D\to\R.$

\begin{lemma}\label{431}
    Let $\Lambda$ be an achronal meridian in $\partial\AS^{2,1}$. Then either $\Lambda$ is the boundary of a lightlike plane, or $S_+(\Lambda)\cap S_-(\Lambda)=\Lambda.$ In the latter case there is an achronal properly embedded surface in $\Omega(\Lambda)$ whose boundary in $\partial\AS^{2,1}$ is $\Lambda$.  
\end{lemma}

\begin{proof}
    Let $f:\pi\D\to\R$ be the function whose graph is $\Lambda.$ We recall from Lemma \ref{417} that $\text{osc}(f)\leq\pi$. If there are points $x_0,x_0^{\prime}$ such that $f(x_0^{\prime})=f(x_0)+\pi$ then combining both Lemma \ref{417} and Lemma \ref{422} we deduce that $\Lambda$ is the boundary of a lightlike lane, and this lightlike plane coincides with $S_+(\Lambda)\cap S_-(\Lambda).$\\
    Assume now that the maximal oscillation of $f$ is smaller than $\pi$, and let us show that $S_+(\Lambda)\cap S_-(\Lambda)=\Lambda.$ By the assumption, if a lightlike geodesic connects $(x_0,f(x_0))$ to $(x_0^{\prime} ,f(x_0^{\prime}))$ then $x_0$ and $x_0^{\prime}$ are not antipodal. But then $x_0,x_0^{\prime}$ are connected by a unique lenght-minimizing geodesic in $\overline{\D}$ for the hemispherical metric, which lies in $\partial\D$. So the lightlike line connecting $(x_0,f(x_0))$ to $(x_0^{\prime},f(x_0^{\prime}))$ is contained in $\partial\AS^{2,1}$. By Lemma \ref{422} we conclude that $S_-(\Lambda)$ and $S_+(\Lambda)$ do not meet in $\AS^{2,1}$ and therefore $S_+(\Lambda)\cap S_-(\Lambda)=\Lambda.$\\
    Finally, in this latter case the function $F=(f^\Lambda_{-+f^\Lambda_+})/2$ is $1-$Lipschitz and defines an achronal properly embedded surface contained in $\Omega(\Lambda)$, whose boundary is $\Lambda$. 
\end{proof}

We remark that in fact for any achronal meridian there is a spacelike surface whose boundary at infinity is $\Lambda$, we will expand on this this in Remark \ref{447}
Now, given a point $x \in \AS^{2,1}$, the Dirichlet domain of $x$ is the region $R_x$ containing $x$ and bounded by two spacelike plane ''dual" to $x$. Namely the planes that we denote (with a slight abuse of notation) $P_x^+$ and $P_x^-$, consisting of points at timelike distance $\pi/2$ in the future (resp. past) along timelike geodesic with initial point $x$.

\begin{proposition}\label{432}
    Let $\Lambda$ be an achronal meridian in $\partial\AS^{2,1}$ different from the boundary of a lightlike plane. Then: 
        \begin{itemize}
            \item A point $x\in\AS^{2,1}$ lies in $\Omega(\Lambda)$ if and only if $\Lambda$ is contained in the interior of the Dirichlet region $R_x$. 
            \item For any $z\in\Lambda$, let $L_-(z)$ and $L_+(z)$ be the two lightlike planes such that $z$ is the past vertex of $L_+(z)$ and the future vertex of $L_-(z)$. Then 
            \[
                \Omega(\Lambda)=\bigcap_{z\in\Lambda}I^+(L_-(z))\cap I^{-}(L_+(z)).
            \] 
            \item The length of the intersection of $\omega(\Lambda)$ with any timelike geodesic of $\AS^{2,1}$ is at most $\pi.$ Moreover, there exist a timelike geodesic whose intersection with $\Omega(\Lambda)$ has length $\pi$ if and only if $\Lambda$ is the boundary at infinity of a spacelike plane. 
        \end{itemize}
\end{proposition}
\begin{proof}
    By Observation \ref{423} a point $x$ lies in $\Omega(\Lambda)$ if and only if it is connected to any point of $\Lambda$ by a spacelike geodesic. The region of points connected to $x$ by a spacelike geodesic has boundary the lightcone from $x,$ whose intersection with $\partial\AS^{2,1}$ coincides with $p_x^\pm\cap \partial\AS^{2,1}.$\\
    Moving on the second item we observe that the region bounded by $L_+(z)$ and $L_-(z)$ contains exaclty points connected to $z$ by a spacelike geodesic. Using the characterization of $\Omega(\Lambda)$ as above, we have the second statement.   \\
    Lastly, if a timelike geodesic $\gamma$ meets $\Omega(\Lambda)$ at a point $x$, then $\Omega(\Lambda)\subset R_x$, so that the lenght of $\gamma\cap\Omega(\Lambda)$ is smaller than the lenght of $\gamma\cap R_x$. But the latter is $\pi.$ Now assume the existance of a geodesic $\gamma$ such that the lenght of $\gamma\cap R_x$ equals $\pi.$ Up to applying an isometry of $\AS^{2,1}$ we may assume that $\gamma$ is vertical in the Poincaré model of $\AS^{2,1}$ amd the mid-point of $\gamma\cap\Omega(\Lambda)$ is $(0,0)$. Thus $(0,-\pi/2)$ and $(0,\pi/2)$ lie on $S_-(\Lambda)$ and $S_+(\Lambda)$ respectively. \\
    Now, again by \ref{423} points of $\Lambda$ are connected to $(0,-\pi/2)$ by a spacelike or lightlike geodesic, hence $s\leq 0$ for all $(\xi,s)\in\Lambda$ Analogously using the point $(0,\pi/2)$ we deduce that $s\geq 0$ for all $(\xi,s)\in\Lambda,$ so that $\Lambda=\partial\D\times\{0\}$.
\end{proof}

Arguing in similar fashion, we obtain that the invisibile domain of an achronal meridian which is not the boundary of a lightlike plane is always contained in a Dirichlet region. 

\begin{proposition}\label{433}
    Given an achronal meridian $\Lambda$ in $\partial\AS^{2,1}$ different from the boundary of a lightlike plane, the invisibile domain $\Omega(\Lambda)$ is contained in a Dirichlet region unless $\Lambda$ is the boundary of a spacelike plane. 
\end{proposition}
\begin{proof}
    We start by defining $a_+=\sup f_+^{\Lambda}$ and $a_-=\inf f_-^{\Lambda}$, and we consider the planes: 
    \[
        Q_{a_+}=\{(x,t)\;|\;t=a_+\} \;\;\text{and}\;\; Q_{a_-}=\{(x,t)\;|\;t=a_-\}
   \]

   in the Poincaré model. Since $\Omega(\Lambda)$ lies in the open region bounded by those planes, it is sufficient to show that $a_+-a_-\leq\pi.$ Assume by contradiction the converse. Notice tht $P_{a_+}$ meets $S_+(\Lambda)$ at some point $p_+=(x_+,a_+)$ and $P_{a_-}$ meets $S_-(\Lambda)$ at some point $p_-=(x_-,a_-)$ where $x_+$ and $x_-$ are points on $\overline{\D}$. For $\epsilon=(a_+ + a_- -\pi )/2$ we can find $x^{\prime}_+$ and $x_-^{\prime} $ in $\D$ such that $p_+^{\prime}=(x^{\prime}_+,a_+-\epsilon)$ and $p_-^{\prime}=(x_-^{\prime},a_- +\epsilon)$ lie in $\Omega(\Lambda)$ (clearly if $x_\pm$ lies in $\D$ we can take $x_\pm^{\prime}=x_\pm)$. As $(a_+-\epsilon)-(a_- - \epsilon )=\pi$, the geodesic segment $\gamma$ joining $p_+^{\prime}$ and $p_-^{\prime}$ is timelike of length $\pi$. Its end-points are in $\text{I}^+(S_-(\Lambda))\cap\text{I}^-(S_+(\Lambda))$, so $\gamma$ is entirely contained in $\Omega(\Lambda)$. As end-points of $\gamma$ are contained in $\Omega(\Lambda).$ $\gamma$ can be extended within $\Omega(\Lambda)$ but this contradicts the third point of Proposition \ref{432}.       
   The third point of Proposition \ref{432} then shows that if $a_+-a_-=\pi$ then $\Lambda$ is the boundary of a spacelike plane. Hence, apart from this case, one has $a_+ - a_-<\pi$, so the closure of $\Omega(\Lambda)$ is contained in a Dirichlet region. 

\end{proof}
    
\begin{observation}
    When $\Lambda$ is the boundary of a spacelike plane $P$, then there are two points $x_-$ and $x_+$ such that $P=P_{x_-}^+=P_{x_+}^-.$ The previous argument shows that $\Omega(\Lambda)$ is the union of all timelike geodesics joining $x_-$ to $x_+$ In this case $S_-(\Lambda$) is the union of the future directed lightlike geodesic rays emanating from $x_-$, whereas $S+(\Lambda)$ is the union of future directed lightlike rays ending at $x_+$.  
\end{observation}

\subsection{Domain of dependence}
We want to define and study properties of Cauchy surfaces and domains of dependence, a priori concept of Lorentzian geometry adapted to the geometry of $\AS^{2,1}$

\begin{definition}
    Given an achronal subset $X$ in a Lorentzian manifold $(M,g)$ the \textit{domain of dependence} of $X$ is the set: 
    \[
        \mathcal{D}(X)=\{p\in M |\;\text{every inextensible causal curve through}\; p\; \text{meets} X \}.
    \]
    We say that $X$ is a \textit{Cauchy surface} of $M$ if $\mathcal{D}(X)=M$. A spacetime $M$ is said \textit{globally hyperbolic} if it admits a Cauchy surface.
\end{definition}

\begin{observation}
    Equivalently the domain of dependence  $D(X)$ of $X$ is the set of points $p\in\A^{2,1}$ such that the plane $p^*,$ the projective dual of $p$ as introduced in \ref{polarsec} is disjoint from $X.$ Domain of dependence are always contained in an affine chart and admit only light-like support planes (for more detail about support planes we will study them in detail in \ref{sectionplanes})
\end{observation}

The theory of globally hyperbolic spacetimes is a well-developed topic in Lorentzian geometry, we will just state the facts that we will need in the following. 

\begin{theorem}\label{442}
    Let M be a globally hyperbolic spacetime. Then:
    \begin{enumerate}
        \item Any two Cauchy surfaces in $M$ are diffeomorphic. 
        \item There exists a submersion $\tau:M\to\R$ whose fibers are Cauchy surfaces.
        \item M is diffeomorphic to $\Sigma\times\R$ where $\Sigma$ is any Cauchy surface in M.
    \end{enumerate}
\end{theorem}

\begin{observation}
    The spacetime $\AS^{2,1}$ is not globally hyperbolic. In fact if $X$ is achronal, it is contained in the graph of a 1-Lipschitz function $f:(\D\cup\partial\D,g_{\mathbb{S}^2})\to\R.$ If $t_0>\text{sup}f$ and $\xi \in \partial\D$, then any lightlike ray with past end-point $(\xi,t_0)$ does not intersect $X.$ 
\end{observation}

\begin{observation}
    As causality notions are invariant under conformal change of metrics, we observe that causal paths in $\AS^{2,1}$ are the graphs of 1-Lipschitz functions from (intervals in) $\R$ to $\D$ with respect to the hemispherical metric in the image. Hence an inextesible causal curve in $\AS^{2,1}$ is either the graph of a global 1-Lipschitz function from $\R$, or it is defined on a proper interval and has endpoint(s) in $\partial\AS^{2,1}.$
\end{observation}

\begin{lemma}\label{445}
    Given an achronal meridian $\Lambda$ in $\partial\AS^{2,1}$ any Cauchy surface in $\Omega(\Lambda)$ is properly embedded with boundary at infinity $\Lambda$. 
\end{lemma}
\begin{proof}
    Let $S$ be a Cauchy surface in $\Omega(\Lambda)$. For every $x\in \D$, the vertical line through $x$ in the Poincaré model meets $\Omega(\Lambda)$, and it interesection with $\Omega(\Lambda)$ must meet $S$ by definition of Cauchy surface. This shows that $S$ is a graph over $\D$, hence properly embedded, and clearly $\partial S=\Lambda.$
\end{proof}

\begin{proposition}\label{446}
    Let $\Lambda$ be an achronal meridian in $\partial\AS^{2,1}$ different from the boundary of a lightlike plane. Let S be a properly embedded achronal surface in $\Omega(\Lambda)$. Then $\mathcal{D}(S)=\Omega(\Lambda)$. In particular $\Omega(\Lambda)$ is a globally hyperbolic spacetime.
\end{proposition}
\begin{proof}
    Let $x$ be any point in $\Omega(\Lambda)$ and take any inextensible causal path through $x$. A priori its future endpoint might be either in $S_+(\Lambda)$ or in $\Lambda$, but by definition of $\Omega(\Lambda$), $x$ cannot be connected by any causal path to $\Lambda$, hence the latter case is excluded. The same argument shows that the past endpoint is in $S_-(\Lambda$.) Since the inextensibile causal path meets both $S_+(\Lambda)$ and $S_-(\Lambda)$, it must meet $S$ by Lemma \ref{445}, hence $\xi\in\mathcal{D}(S)$. \\
    Conversely, consider a $x$ that is not in $\Omega(\Lambda)$, the one can find a causal path joining $x$ to $\Lambda$, which is necessarily inextensibile. Hence $x$ is not in $\mathcal{D}(S)$.  
\end{proof}

\begin{observation}\label{447}
    As a direct consequence of Theorem \ref{442} and Proposition \ref{446} we have that $\Lambda$ is the boundary of a spacelike surface in $\Omega(\Lambda)$, namely a Cauchy surface in $\Omega(\Lambda)$. By Lemma \ref{445}, the surface is properly embedded, hence the graph of a global 1-Lipschitz function. This shows that any proper achronal meridian $\Lambda$ is the boundary at infinity of a properly embedded spacelike surface, we remark that we have improved the statement of Lemma \ref{431}.
\end{observation}

The most remarkable property of domain of dependence of a properly embedded surface in $\AS^{2,1}$, and it is a direct consequence of Proposition \ref{446}, is that it only depends on the boundary at infinity. In detail:

\begin{corollary}\label{448}
    If $S$ and $S^{\prime}$ are properly embedded spacelike surfaces in $\AS^{2,1}$ , then $\mathcal{D}(S)=\mathcal{D}(S^{\prime})$ if and only if $\partial S=\partial S^{\prime}$.
\end{corollary}

\subsection{Properly achronal sets in $\A^{2,1}$}
For our interest will be important to consider the model $\A^{2,1}$. As $\A^{2,1}$ contains closed timelike lines, it does not contain any achronal subset. However if $P$ is a spacelike plane in $\A^{2,1}$ then $\A^{2,1}\setminus P$ does not contain closed causal curves. Indeed it is simply connected, so it admits an isometric embedding into $\AS^{2,1}$, given by a section of the covering map $\AS^{2,1}\to\A^{2,1},$ and whose image is a Dirichlet region.
\begin{definition}
    A subset $X$ of $\A^{2,1}\cup\partial\A^{2,1}$ is a \textit{proper achronal subset} if there exists a spacelike plane $P$ such that $X$ is contained in $\A^{2,1}\cup\partial\A^{2,1}\setminus\overline{P}$ and is achronal as a subset of $\A^{2,1}\cup\partial\A^{2,1}\setminus\overline{P}.$  
\end{definition}

It follows from the definition that if $X$ is a proper achronal subset of $\overline{\A^{2,1}}$, then it admits a section to $\AS^{2,1}\cup\partial\AS^{2,1}$, and the image remains achronal after the lifting. Conversely, let $\tilde{X}$ be an achronal subset of $\AS^{2,1}$ different from a lightlike plane, then it is contained in a Dirichlet region as a consequence of Lemma \ref{417}, and the fact that any achronal subset of $\AS^{2,1}$ is contained in properly embedded one. As Dirichlet regions are projected in $\AS^{2,1}$ to the complement of a spacelike plane, the image of $\tilde{X}$ to $A^{2.1}$ is a proper achronal subset. The following lemma will be key in our \textit{path} to prove the earthquake theorem:\\

\begin{lemma}
    Let $\varphi:\R\text{P}^1\to\R\text{P}^1$ be an orientation preserving homeomorphism. Then the graph of $\varphi$, say $\Gamma_\varphi\subset \R\text{P}^1\times\R\text{P}^1\simeq\partial\A^{2,1}$ is a proper achronal subset and any lift $\tilde{\Gamma_\varphi}$ is an achronal meridian in $\partial\AS^{2,1}.$  
\end{lemma}
\begin{proof}
    We start by proving that $\Lambda_\varphi$ is locally achronal. Consider $U$ and $V$ intervals around $x$ and $@\varphi(x)$ and $\theta_1$ and $\theta_2$ are positive coordinates on $U$ and $V$ respectively, then timelike curves $\gamma(t)=(\gamma_1(t),\gamma_2(t))$ in $U\times V$ are characterized by the porpery that $\theta_1^{\prime} (t)\theta_2^{\prime} (t)<0,$ where $\theta_i(t)=\theta_i(\gamma_i(t)).$ \\
    In particular points on $\Lambda_\varphi\cap U\times V$ are not related by a timelike curve contained in $U\times V$, by the assumption that $\varphi$ is orientation-preserving. \\
    Let us prove the existance of a spacelike plane $P$ such that $\overline{P}\cap \Lambda_\varphi=\emptyset.$ Let us consider the identification $\S=\R\cup\{\infty\}$, and take $\varphi_0\in\PSL(2,\R)$ so that $\varphi_0^{-1}\varphi(0)=1, \varphi_0^{-1}\varphi(1)=\infty$ and $\varphi_0^{-1}\varphi(\infty)=0$ It follows that $\varphi_0^{-1}\varphi$ sends the intervals $(\infty,0), (0,1)$ and $(1,\infty$) respectively to $(0,1),(1,\infty),(\infty,0)$. Thus $\varphi_0^{-1}\varphi$ has no fixed points, that is, the graph of $\varphi$ does not intersect the graph of $\varphi_0$, which is the asymptotic boundary of the spacelike plane $P_{\varphi_0}.$\\
    Let us consider now the lift $\widetilde{\Lambda}_\varphi$ to the boundary of $\AS^{2,1}$. As $\Lambda_\varphi$ is contained in a simply connected region of $\overline{\A^{2,1}}$, its lift is a closed locally achronal curve contained in $\partial\AS^{2,1}$. In particular the projection $\widetilde{\Lambda}_\varphi\to\partial\D$ is locally injective. As $\widetilde{\Lambda}_\varphi$ is compact, the map is a covering. On the other hand, since $\Lambda_\varphi$ is homotopic to the boundary of a plane in $\partial\A^{2.1}$, it turns out that $\widetilde{\Lambda}_\varphi$ is homotopic to $\partial\D$ in $\partial\AS^{2,1}$, so that the projections $\widetilde{\Lambda}_\varphi\to\partial\D$ is bijective. It follows that $\widetilde{\Lambda}_\varphi$ is achronal.  
\end{proof}

What we have shown until now for achronal sets in $\AS^{2,1}$ can be rephrased for proper achronal sets of $\A^{2,1}$. For example, any proper achronal set $X$ can be extended to a properly embedded proper achronal surface and there are two extremal extensions, just as in Lemma \ref{422}\\
We would like to focus now on proper achronal meridians in $\partial\A^{2,1}$. They lift to achronal meridians in $\partial\AS^{2,1}$ different from the boundary of lightlike planes. Indeed the boundary of a lightlike plane is not contained in a Dirichlet region. Conversely any achronal meridian on $\partial\AS^{2,1}$ different from the boundary of a lightlike plane projects to an achronal meridian of $\A^{2,1}$.

\begin{proposition}\label{453}
    Let $\Lambda$ be a proper achronal meridian in $\partial\A^{2,1}$ and denote by $\widetilde{\Lambda}$ any lift to the universal covering. Then the universal covering map of $\A^{2,1}$ maps $\Omega(\widetilde{\Lambda})$ injectively to the domain: 
    \[
        \Omega(\Lambda)\coloneqq \{x\in\A^{2,1}\;|\;P_x\cap\Lambda=\emptyset\}.
    \]
\end{proposition}
    
\begin{proof}
Consider the covering $p:\AS^{2,1}\to \A^{2,1}$, by Propostion \ref{433} the invisible domain $\Omega(\widetilde{\Lambda})$ is contained in a Dirichlet region $R_{\widetilde{x}}$, hence the restricion of $p$ to $\Omega(\widetilde{\Lambda})$ is injective and its image is contained in $p(R_{\widetilde{x}}),$ namely the complement in $\overline{\A^{2,1}}$ of the spacelike plane $P_x$ dual to $x=p(\widetilde{x})$. Moreover, by the firs item of Proposition \ref{432}, one can actually pick for $\widetilde{x}$ any point in $\Omega(\widetilde{\Lambda})$, which shows that the image $p(\Omega(\widetilde{\Lambda}))$ is contained in $\Omega(\Lambda)$ defined as in the proposition. \\ 
For the converse inclusion, let $x\in\A^{2.1}$ be a point whose dual plane $P_x$ does not meet $\Lambda$. The preimage $p^{-1}(P_x)$ is a countable disjoint union of planes which disconnect $\AS^{2,1}\cup \partial\AS^{2,1}$ in a disjoint union of Dirichlet regions centered at preimages of $x$. The lifr $\widetilde{\Lambda}$ is then contained in exactly one such region, say $R_{\widetilde{x}}$. By the first point of \ref{432} $\widetilde{x}\in\Omega(\widetilde{\Lambda})$ which implies that   $x=p(\widetilde{x})$ lies in $p(\Omega(\widetilde{\Lambda}))$. 
\end{proof}


When $\Lambda$ is the graph of an orientation-preserving homomorphism $\varphi:\S\to\S$, there is a fairly simple characterization of $\Omega(\Lambda)$ exploiting the identification $\A^{2,1}\simeq \PSL$.


\begin{corollary}\label{454}
    Let $\varphi:\S\to\S$ be an orientation-preserving homeomorphism. Then $x\in\A^{2,1}$ lies in $\Omega(\Lambda_\varphi)$ if and only if $x\circ\varphi$ has no fixed point as a homeomorphism of $\S$.
\end{corollary}
\begin{proof}
    The dual plane of $x,$ viewed as an element of $\PSL$, meets $\partial\A^{2,1}$ along the graph of $x^{-1}$, namely $\Lambda_{x^{-1}}$.\\
    With this remark in hand, we have that $x\in\Omega(\Lambda_\varphi)$ if and only if $\Lambda_{x^{-1}}\cap \Lambda_\varphi=\emptyset$, that is, if and only if $x^{-1}\circ\varphi$ has no fixed point on $\S$.

\end{proof}

\begin{proposition}
    Let $\sigma: S\to\A^{2,1}$ be a proper spacelike immersion. Then 
    \begin{itemize}
        \item $\sigma$ is a proper embedding.
        \item $\sigma$ lifts to a proper embedding $\tilde{\sigma}:S\to\AS^{2,1}$.
        \item The boundary at infinity of $\sigma(S)$ is a proper achronal meridian $\Lambda$ in $\partial\A^{2,1}$.
        \item $\mathcal{D}(\sigma(S))=\Omega(\Lambda).$
    \end{itemize}
\end{proposition}

\begin{proof}

    Denote by $\widehat{S}$ the covering of $S$ admitting a lift $\widehat{\sigma}:\widehat{S}\to\H^{2.1}$. In general either $\widehat{S}=S$ or it is a $2:1$ covering. Since the covering is finite, $\widehat{\sigma}$ is a proper immersion. Let us consider the identification $\H^2\times\mathbb{S^1}$ that follows from \refeq{ogcover}. The induced projection: $\text{pr}:H^{2,1}\to\H^2$ is a proper fibration with timelike fibers. In particular $\widehat{\sigma}$ is trasverse to the fibers of $\text{pr}.$ It follows that the composition $\text{pr}\circ\sigma:\widehat{S}\to\H^2$ is a proper local diffeomorphism, hence a covering map. Since $H^2$ is simply connected, we deduce that the covering is actually a homeomorphism, $\widehat{\sigma}$ is an embedding, and $\widehat{S}$ is homeomorphic to the hyperbolic plane. \\
    In particular we can lift $\widehat{\sigma}$ to the universal covering, say $\widehat{\sigma}:\widehat{S}\to\AS^{2,1},$ which is still a proper spacelike embedding $\widehat{S}\to \AS^{2,1}$. By Lemma \ref{413} and Lemma \ref{415} we know that the image is an achronal meridian, and it is contained in a Dirichlet domain by Lemma \ref{417}. It follows that $\widehat{\sigma}(\widehat{S})$ is contained in a Dirichlet domain of the covering map $\H^{2,1}\to\A^{2,1}$, on which we know that the covering map is injective. In particular $\sigma$ is injective, hence $S=\widehat{S}$ as desired.  
\end{proof}

\begin{observation}
    In the proof of the last proposition, once we proved that $\hat{S}$ is homeomorphic to $\R^2,$ we could have inferred immediately that $\hat{S}=S$ as $\Z/2\Z$ cannot act freely on $\R^2$ by diffeomorphisms. 
\end{observation}

We therefore have an analogue version of Lemma \ref{448} in $\A^{2,1}$.

\begin{corollary}
    If $S$ and $S^{\prime}$ are properly embedded spacelike surfaces in $\A^{2,1}$, then $\mathcal{D}(S)=\mathcal{D}(S^{\prime})$ if and only if $\partial S=\partial S^{\prime}$. 
\end{corollary}

\subsection{Convexity properties.}

Let $\Lambda$ be a proper achronal meridian in $\partial\A^{2,1}$. We would like to investigate convexity properties of $\Omega(\Lambda)$. We briefly recall that given $X \subset \R\text{P}^n$ is convex if it contained in an affine chart, and it is convex in the affine chart. The notion does not depend on the affine chart containing $X$. We would say that it is a proper convex set if it compactly contained in an affine chart. 

\begin{proposition}\label{461}
    Given a proper achronal meridian $\Lambda$ in $\partial\A^{2,1}$, the invisible domain $\Omega(\Lambda)$ is convex. If $\Lambda$ is different from the boundary of a spacelike plane then $\Omega(\Lambda)$ is a proper convex set.
\end{proposition}

\begin{proof}
    By Proposition \ref{433} there exists a spacelike plane $P$ such that $\Omega(\Lambda)$ is contained in the affine chart $V$ of $\R\text{P}^3$ obtained by removing the projective plane containing $P$. The domain $\A^{2,1}\cap V=\A^{2,1}\setminus P$ is isometric to a Dirichlet region $R$ of $\AS^{2,1}$, by an isometry that sends $\Lambda$ to a lifting $\widetilde{\Lambda}$ and $\Omega(\Lambda)$ to $\Omega(\widetilde{\Lambda})$. By te second point of Proposition \ref{432} we have 
    \[
        \Omega(\widetilde{\Lambda})=\bigcap_{\widetilde{z}\in\widetilde{\Lambda}}I^+(L_-(\widetilde{z}))\cap I^-(L_+(\widetilde{z})).
    \]

    Now if $\widetilde{z}$ projects to $z$, then the images of $L_-(\widetilde{z})$ and $L_+(\widetilde{z})$ in $V$ are the two components of $L(z)\cap \A^{2,1}$, where $L(z)$ is the affine tangent plane of $\partial\A^{2,1}\cap V$ at $z$. It turns out that te image of the region $I^+(L_-(\widetilde{z}))\cap I^-(L_+(\widetilde{z}))$ is the interesection of $\A^{2,1}$ with the half-space $U(z)$ bounded by $L(z)$ and whose closure contains $\Lambda$. This shows: 
\[
    \Omega(\Lambda)=\A^{2,1}\bigcap_{z\in\Lambda}U(z).
\] 

Actually we claim that:
\[
    \Omega(\Lambda)=\bigcap_{z\in\Lambda}U(z)\subset \A^{2,1},
\]
and this will conclude. As $\bigcap_{z\in\Lambda}U(z)$ is connected and meets $@A^{2,1}$, to show that is is contained in $@A^{2,1}$ it sufficies to show that it does not meet the boundary. For any $w\in\partial\A^{2,1}$ let us consider the leaf of the left ruling through $w$, which intersects $\Lambda$ at a point $z$. It turns out that $L(z)$ contains the leaf of the left ruling through $z$, hence $w\notin U(z)$. \\
Now assume that $\Lambda$ is not the boundary of a spacelike plane. Then by Proposition \ref{433} on the universal covering the compact set $\Omega(\widetilde{\Lambda})\cup S_+(\widetilde{\Lambda})\cup S_-(\widetilde{\Lambda})$ is contained in a Dirichlet domain, so its image is a compact set contained in an affine chart.
\end{proof}

As a consequence we have that $\Lambda$ is contained in an affine chart whose complement in $\R\text{P}^3$ is a projective plane containing a spacelike plane of $\A^{2,1}$ Hence it makes sense to give the following definition:

\begin{definition}
    Given a proper achronal meridian $\Lambda$ in $\partial\A^{2,1}$, we define $\mathcal{C}(\Lambda)$ to be the convex hull of $\Lambda$, which can be taken in an affine chart containing $\Lambda$
\end{definition}

What we have implicitly proved is that given $\Lambda$ an achronal meridian in $\partial\A^{2,1}$, then $\mathcal{C}(\Lambda)$ is contained in $\A^{2,1}$, which is not immediately obvious as $\A^{2,1}$ is not convex in $\R\text{P}^3$. \\

\begin{example}\label{43}
    Consider $\sigma\in\PSL$, the convex hull of the graph $\Lambda_\sigma$ is the closure of the totally geodesic spacelike plane $P_{\sigma^{-1}}$ in $\A^{2.1}$. In particular, following Lemma \ref{32}, we have that the boundary at infinity of $P_{\sigma^{-1}}$ equals $\Lambda_\sigma$, and moreover $P_{\sigma^{-1}}$ is convex, since spacelike geodesics of $\A^{2,1}$ (which are the intersection of two transverse spacelike planes) are lines in an affine chart, and any two points in $\partial\H^2$ are connected by a geodesic. Hence $P_{\sigma^{-1}}$ is clearly the smallest convex set containing $\Lambda_\sigma$.  
\end{example}

\begin{observation} Since $\Omega(\Lambda)$ is convex, $\mathcal{C}(\Lambda)$ is contained in $\Omega(\Lambda)$. Moreover, if $K$ is any convex set contained in $\overline{\A^{2,1}}$ and containing $\Lambda$, then $\mathcal{C}(\Lambda)\subset K\subset \overline{\Omega(\Lambda)}.$\\
    To see this, let $V$ be an affine chart such that $\Lambda\subset V$ is obtained removing a spacelike projective plane. Now, if $z\in \Lambda$ then for any $x\in\A^{2,1}\cap V$ the segment connecting $z$ and $x$ in $V$ is contained in $\A^{2,1}$ if and only if $x\in U(z)$, the half-space containing $\Lambda$ and bounded by the tangent space of $\Lambda$ at $z,$ as in the proof of Proposition \ref{461}. \\
    It follows, from the characterization of $\Omega(\Lambda)$ as the intersection of the $U(z)$ given in Proposition \ref{461}, that if $x$ is not in $\overline{\Omega(\Lambda)}$ it cannot be in $K$. Hence $\overline{\Omega(\Lambda)}$ is the biggest convex set of $\A^{2,1}$ containing $\Lambda$.
\end{observation}

Suppose now that $\Lambda$ is not the boundary of a spacelike plane. Now the topological frontiers in $\R\text{P}^3$ of $\Omega(\Lambda)$ and of $\mathcal{C}(\Lambda)$ are Lipschitz surfaces homeomorphic to a sphere. This sphere is disconnected by $\Lambda$ in two regions, homeomorphic to disks, which form the boundary of $\Omega(\Lambda)$ and of $\mathcal{C}(\Lambda)$ in $\A^{2,1}$. For $\Omega(\Lambda)$ those components are the image of $S_\pm(\widetilde{\Lambda})$ and will be denoted by $S_\pm(\Lambda$). \\
For $\mathcal{C}(\Lambda),$ let be $\mathcal{C}(\widetilde{\Lambda})$ a lifting, which is necessarily contained in a Dirichlet region, say $R$. Let $P$ be a support plane for $\mathcal{C}(\Lambda)$, which is necessarily either spacelike or lightlike, and let $\widetilde{P}$ be its lift which touches $\mathcal{C}(\widetilde{\Lambda}).$ Hence either $\widetilde{\Lambda}$ is in $I^+(\widetilde{P})\cup\widetilde{P}$ or in $I^-(\widetilde{P})\cup\widetilde{P}$. This permits to distinguish the componets of $\partial\mathcal{C}(\Lambda)\setminus\Lambda:$ the \textit{past boundary component} $\partial\mathcal{C}_-(\Lambda)$ has the property that $\widetilde{\Lambda}$ is contained in $I^+(\widetilde{P})\cup \widetilde{P}$ for all support planes which touch $\partial_-\mathcal{C}(\Lambda)$, and analogously for the \textit{future boundary component} $\partial_+\mathcal{C}(\Lambda)$.\\
We want to show a kind of duality between boundary components $\partial_\pm\mathcal{C}(\Lambda)$ and $S_\pm(\Lambda)$. 

\begin{proposition}\label{464}
Let $\Lambda$ be a proper achronal meridian in $\A^{2,1}$, $x\in\A^{2,1}$ and let's denote by $P_x$ the dual plane to $x$. Then: 
\begin{itemize}
    \item $x\in\Omega(\Lambda)$ if and only if $P_x\cap\mathcal{C}(\Lambda)=\emptyset$. 
    \item $x\in\mathcal{C}(\Lambda)$ if and only if $P_x\cap \Omega(\Lambda)=\emptyset$.
\end{itemize} 
In particular if $\Lambda$ is not the boundary of a spacelike plane, then: 
\begin{itemize}
    \item $x\in\partial_\pm\Omega(\Lambda)$ if and only if $P_x$ is a support plane for $\partial_\mp\mathcal{C}(\Lambda).$
    \item $x\in\partial_\pm\CF$ if and only if $P_x$ is a support plane for $S_\mp(\Lambda)$.
\end{itemize}
\end{proposition}
\begin{proof}
From Proposition \ref{453}, points in $\Omega(\Lambda)$ are dual to planes disjoint from $\Lambda$, which are precisely those which do not intersect $@\mathcal{C}(\Lambda)$, by the definition of convex hull. For the second statement, fix $x$ and observe that $z\in P_x$ if and only if $x\in P_z$. Hence there exists a point $z$ in the intersection $P_x\cap \Omega(\Lambda)$ if and only if $x$ is in a plane $P_z$ which is disjoint from $\Lambda$, namely when $x$ is not in $\mathcal{C}(\Lambda)$. \\
As a consequence $\partial\mathcal{C}(\Lambda)$ consists of points dual to support planes of $\Omega(\Lambda)$. Take a support plane $P_x$ of $S_+(\Lambda)$ (hence dual to a point $x$) which meets $S_+(\Lambda)$ at $z$. If $\widetilde{z}$ denotes the corresponding point of $S_+(\widetilde{\Lambda}),$ then $\widetilde{\Lambda}\subset I^+(P_{\widetilde{z}}^-),$ and $P_{\widetilde{z}}^-\cap\widetilde{\Lambda}=\emptyset.$ Thus $P_z$, which is the projection of $P_{\widetilde{z}}^-$, is a support plane of $\mathcal{C}(\Lambda)$ touching the past boundary. As $x\in P_z$ we conclude that $x$ lies in the past boundary. Similary points of the future boundary of $\mathcal{C}(\Lambda)$ correspond to support planes for $S_-(\Lambda)$.
\end{proof}

\begin{observation}\label{465}
    It may happen that a boundary component of $\mathcal{C}(\Lambda)$ meets the boundary of $\Omega(\Lambda)$. This exactly happens when the curve $\Lambda$ contains a \textit{sawtooth}, namely two consecutive lightlike segments in $\A^{2,1},$ one past directed and the other future-directed. In this case the lighltike plane $L(z)$ tangent to $\partial\A^{2,1}$ at the vertex $z$ of the sawtooth contains a lightlike traingle contained in $L(z)$. This is however not contained in $\Omega(\Lambda)$. If the curve $\Lambda$ does nto contain any sawtooth, then $@\mathcal{\Lambda}\setminus\Lambda$ is entirely contained in $\Omega(\Lambda)$.\\
    The fundamental example is given in \textcolor{red}{inserire figura 8}, where the yellow refion represents at the same time the convex hull of the proper achronal meridian $\Lambda$ in $\partial\A^{2,1}$ composed of four lightlike segments, two past-directed and two future-directed, and the closure of $\Omega(\Lambda).$ \textcolor{red}{Vedi anche remark 8.1.3 e la figura 14.} 
\end{observation}

\begin{proposition}\label{466}
The past and future boundary components of $\mathcal{C}(\Lambda)$ are achronal surfaces.
\end{proposition}
\begin{proof}
    We will show it for $\partial_+\mathcal{C}(\Lambda)$. Take $x,y\in\partial_+\mathcal{C}(\Lambda)$ and consider the segment joining $x$ to $y$ in an affine chart containing $\Lambda$. If this segment is timelike then the dual planes $P_x$ and $P_y$ would be disjoint. Then (up to switching the roles of $x$ and $y$) we may assume that, in the universal cover $P_{\widetilde{x}}^1\subset I^+(P_{\widetilde{y}}^1)$, where $\widetilde{x}$ and $\widetilde{y}$ are the lifting of $x$ and $y$ in the same Dirichlet region mapping to the fixed affine chart. But then $S_+(\widetilde{\Lambda})$ would be contained in $\text{I}^-(P_{\widetilde{y}}^1)$ and could not meet $P_{\widetilde{x}}^\$$ thus a contradiction to Proposition \ref{464}
\end{proof}

\begin{observation}\label{467}
The past and future boundary components of $@\mathcal{C}(\Lambda)$ are not smooth, but only Lipschitz surfaces. Indeed the complement of $\Lambda$ and of the lightlike triangles (as described in \ref{465}) is locally connected by acausal Lipschitz arcs, and one can define a pseudo-distance, that in fact turns out to be a distance and makes $@\mathcal{C}(\Lambda)$ locally isometric to the hyperbolic plane. \\
The situation is very similar to the counterpart in hyperbolic three-space. The intersection of a spacelike support plane with $\mathcal{C}(\Lambda)$ is either a geodesic or a straight convex subset of $\H^2$, i.e. a subset bounded by geodesics. Thus $\partial\mathcal{C}(\Lambda)\setminus\Lambda$ is intrinsically a hyperbolic surface pleated along a measured geodesic lamination. A remarkable difference with respect to the hyperbolic case \textcolor{red}{si ma quale} is that in general those surfaces may be not complete, but they are always isometric to straight convex subset of $\H^2$. \textcolor{red}{Per espandere su questo consigliano Multi-black holes and earthquakes on Riemann surfaces with boundaries.}

\end{observation}

\section{Support planes.}\label{sectionplanes}
We still need to borrow some notions and notations from convex analysis. Given a convex body $K$ in affine space of dimension three, a \textit{support plane} of $K$ is an affine plane $Q$ such that $K$ is contained in a closed half-space bounded by $Q$, and $\partial K\cap Q\neq\emptyset.$ If $p$ is in the intersection of $K$ with $Q$ we will say that $Q$ is a \textit{support plane} at $p$. As a consequence of the Hahn-Banach theorem there exixst a support plane at every point $p\in\partial K.$ \\
We will adopt the terminology to the Anti-de Sitter setting, given a convex hull $\mathcal{C}(f)$ in $\A^{2,1}$, we say that a totally geodesic plane $P$ is a support plane of $\mathcal{C}(f)$ (at $p\in\partial\mathcal{C}(f))$ is $p\in\mathcal{C}(f)\cap \overline{P}\subset\overline{\A^{2,1}}$ and, in an affine chart containing $\text{graph}(f)$, $\mathcal{C}(f)$ lies in a closed half-space bounded by the affine plane that contains $P$. Even this definition does not depend on the choice of the affine chart as long this one contains $\text{graph}(f)$.\\ 

\textcolor{blue}{He goes for a remark that we might not need}\\

Recall \textcolor{orange}{from where?} that if $X$ is a set, $\mathcal{C}(X)$ its convex hull and $Q$ an affine support plane for $\mathcal{C}(X)$, then $Q\cap\mathcal{C}(X)=\mathcal{C}(Q\cap X).$ Applying this identity in our setting, we get that for any totally geodesic support plane $P$: 

\begin{equation}\label{suplane}
    P\cap\mathcal{C}(f)=\mathcal{C}(\partial P\cap\text{graph}(f)).
\end{equation}
    
We can characterize support planes by considering intersection with the convex hull as follows: 

\begin{proposition}\label{supportinho}
    Let $f:\S\to\S$ be an orientation-preserving homeomorphism, and let $P$ be a support plane of $\mathcal{C}(f)$ at a point $p\in\partial\mathcal{C}(f).$ Then:
    \begin{itemize}
        \item If $p\in\A^{2,1}$, then $P$ is a spacelike plane. 
        \item If $p\in\partial\A^{2,1}$ then $P$ is either spacelike or lightlike.
    \end{itemize}
\end{proposition}
\begin{proof}
    The key observation is that if $P$ is a spacelike plane, then $\partial P$ and $\text{graph}(f)=\mathcal{C}(f)\cap\partial\A^{2,1}$ do not intersect transversely. From \textcolor{blue}{add reference}, if $P$ is timelike then $\partial P$ is the graph of an orientating-reversing homeomorphism of $\S$, hence it intersects $\text{graph}(f)$ transversally. From Lemma \textcolor{blue}{add reference}, if $P$ is lightlike, then its boundary is the union of the two circles: $\{x\}\times\S$ and $\S\times{y}$. So if $p\in\partial P\cap\text{graph}(f)$ and $p$ is not the point $p_0=(x,y),$ then $\partial P$ and $\text{graph}(f)$ intersect transversally. Hence the only case where $P$ can be a lightlike support plane is when it intersects $\text{graph}(f)$ only at $p_0.$ We are left with the task to show that $P\cap\mathcal{C}(f)=\{p_0\}$ and it does not contain any point of $\A^{2,1}.$ By contradiction suppose the existance of a $q\in P\cap\mathcal{C}(f)$ with $q$ different from $p_0,$ then by \refeq{suplane} $\partial P\cap\text{graph}(f)$ would also contain another point different from $p_0$, a contradiction.
\end{proof}

Now, given a spacelike support plane $P$ of $\mathcal{C}(f)$ at a point $p$, we say that $P$ is a \textit{future} (resp. \textit{past}) \todo{Forse la definizione del survey è migliore?} support plane if in a small simply connected neighbourhood of
$p\in\overline{\A^{2,1}}, \mathcal{C}(f)$ is contained in the closure of the connected components of $U\setminus P$ which is the past (resp. future) of $P$. This means that there exists a future-oriented (resp past-oriented) timelike curves leaving $\mathcal{C}(f)\cap U$ and reaching $P\cap U$. \\
We observe that $\mathcal{C}(f)$ cannot have both a future and past support plane at $p$ unless $\CF$ has empty interior, a situation that happens exactly when $f$ is an element of $\PSL$, the example treated previously.\\
We can now state the following related to convergence of support planes: 
\begin{lemma}\label{49}
    Let $f:\S\to\S$ be an orientation-preserving homeomorphism which is not in $\PSL$, $p_n$ a sequence of points in $\partial\CF$, and $P_{\gamma_n}$ a sequence of future (resp. past) spacelike supports planes at $p_n$, for $\gamma_n\in\PSL$. Up to extracting a subsequence, we can assume $p_n\to p$ and $P_{\gamma_n}\to P.$ Then: 
    \begin{itemize}
        \item If $p\in\A^{2,1}$, then $P=P_\gamma$ is a future (resp. past) support plane of $\CF$, for $\gamma_n\to\gamma\in\PSL$.
        \item If $p\in\partial\A^{2,1}$, then either $P$ is a lighlike plane whose boundary is the union of two circle meeting at $p$ or the previous point holds.        
    \end{itemize}
\end{lemma}
\begin{proof}
    Diaf-Seppi
\end{proof}
\begin{corollary}\label{nametag}
    Let $f:\S\to\S$ be an orientation-preserving homeomorphism which is not in $\PSL$. Then $\partial\CF$ is the disjoint union of $\text{graph}(f)=\CF\cap\partial\A^{2,1}$ and of two topological discs, of which one admits future support plane, and the other only admits past support plane. 
\end{corollary}
\begin{proof}
    Diaf-Seppi
\end{proof}

According to Corollary \ref{nametag} we will call the connected component of $\partial\CF\setminus\text{graph}(f)$ that only admits future support planes the \textit{future boundary component}, and denote it by $\partial_+\CF$. Similarly, we will call the connected component of $\partial\CF\setminus\text{graph}(f)$ that only admits past support planes the \textit{past boundary component}, and denote it by $\partial_-\CF$.\\

\section{Globally Hyperbolic three-manifolds}
We want now to focus our attention on maximal globally hyperbolic (MGH) Anti-de Sitter spacetimes containing a compact Cauchy surfaces of genus $r$ (we will with a slight abuse of notation say that the spacetime has genus $r$.) We will show that there are no such manifolds when $r=0$. Furthermore, we will then give a rapid overview over the case $r=1$ and the focus on the most interesting case, $r\geq 2$, that will lead to a complete classification.

\subsection{General facts.} We will now state some general that will be useful in our classification.

\begin{lemma}\label{properembedding}
    Let $\sigma:S\to\A^{2,1}$ be a spacelike immersion. If $\sigma^*(g_{\A^{2,1}})$ is a complete Riemannian metric, then $\sigma$ is a proper embedding and $S$ is diffeomorphic to $\R^2$
\end{lemma}
\begin{proof}
    
\end{proof}

As an immediate consequence there cannot be any globally hyperbolic spacetime with genus $0$. In fact, suppose such a spacetime exists and denote by $\Sigma$ the Cauchy surfaces diffeomorphic to $\mathbb{S}^2$, the developing map restricted to $\Sigma$ would be a spacelike immersion, and the pull-back metric would be complete by compactness. But this contradicts \ref{properembedding}. Hence: 

\begin{corollary}
    There exists no globally hyperbolic Anti-de Sitter spacetime of genus 0. 
\end{corollary}

The following is a fundamental result on the structure of globally hyperbolic AdS spacetimes.

\begin{proposition}\label{holorep}
    Let $M$ be a globally hyperbolic Anti-de Sitter spacetimes of genus $r\geq 1$. Then: 
    \begin{enumerate}
        \item The developing map $\text{dev}:\tilde{M}  \to  \A^{2,1}$ is injective. 
        \item If $\Sigma$ is a Cauchy surface of $M,$ then the image of dev is contained in $\Omega(\Lambda)$, where $\Lambda$ is boundary at infinity of $\text{dev}(\tilde{\Sigma}).$
        \item If $\rho:\pi_1(M)\to\text{Isom}(\A^{2,1,})$ is the holonomy representation, $\rho(\pi_1(M))$ acts freely and properly discontinuosly on $\Omega(\Lambda)$, and $\Omega(\Lambda)$ is a globally hyperbolic spacetime containing $M$. 
    \end{enumerate}
\end{proposition}

\begin{proof}
    
\end{proof}

A remarkable difference between Lorentianz and Riemannian geometry is the that in Lorentianz geometry completeness is a very strong assumption, \textcolor{red}{reference} and in fact interesting classification results are obtained removing such condition. However, it is necessary to impose some maximality condition to compensate for non-completeness. Among several  approaches, one of the most common is the classification of a maximal globally hyperbolic spacetimes. Following \cite{bonsanteseppi} we restrict to a less general setting, but one could give similar definitions in the lager class of Einstein spacetimes.

\begin{definition}
    A globally hyperbolic Anti-de Sitter manifold $(M,g)$ is \textit{maximal} if any isometric embedding of $(M,g)$ into a globally hyperbolic Anti-de Sitter manifold $(M^{\prime},g^{\prime})$ which sends a Cauchy surface of $M$ to a Cauchy surface of $M^{\prime}$ is surjective.
\end{definition}

In eyesight of this definition and as a direct consequence of propostion \ref{holorep} we have:

\begin{corollary}
An Anti-de Sitter globally hyperbolic spacetime $M$ is maximal is and only if $\tilde{M}$ is isometric to the invisibile domain of a proper achronal meridian in $\partial\A^{2,1}.$ \textcolor{red}{importante da capire credo}
\end{corollary}


%genere uno 

\subsection{Examples of genus $r\geq 2$}
Let's $\Sigma_r$ be an oriented surface of genus $r\geq 2.$ We recall the definiton of Fuchsian representation. 

\begin{definition}
    A representation $\rho:\pi_1(S)\to \text{PSL}(2,\R)$ is called positive \textit{Fuchsian} if there is a $\rho$-equivariant orientation preserving homeomorphism $\delta:\widetilde{\Sigma}_r\to\H^2$.
\end{definition}

The definition is invariant under conjugation in $\text{PSL}(2,\R)\simeq\text{Isom}_0(\H^2),$ but not under conjugation in $\text{Isom}(\H^2).$ It is a result of Goldman \textcolor{red}{Reference}, that a representation $\rho$ is positive Fuchsian if and only if the associated flat $\R\text{P}^1$ bundle $E_\rho$ constructed as the quotient of $\widetilde{\Sigma}_r\times\R\text{P}^1$ by the diagonal action of $\pi(S)$ (via deck transformation and the representation) has Euler class $2-2r.$ An equivalent condition is the existence of an orientation-preserving fiber bundle isomorphism between $E_\rho$ and the unit tangent bundle of $\Sigma_r.$ \\   

Another useful tool for the classification of genus $r\geq 2$ is the following fact in Teichm\"uller theory \textcolor{red}{reference}:

\begin{lemma}
    Given two positive Fuchsian representation $\rho_l,\rho_r:\pi_1(\Sigma_r)\to\text{PSL}(2,\R)$, any $(\rho_l,\rho_r)$ equivariant orientation-preserving homeomorphism of $\H^2,$ extends continuously to an orientation-preserving homeomorphism of $\H^2\cup\R\text{P}^1.$ Moreover, its extension $\varphi:\R\text{P}^1\to\R\text{P}^1$ is the unique $(\rho_l,\rho_r)-$equivariant orientation homeomorphism of $\R\text{P}^1$.     
\end{lemma}

Here by ($\rho_l,\rho_r$)-equivariant we mean that $\varphi\circ\rho_l=\rho_r\circ\varphi.$\\
Now let $\rho_l,\rho_r$ be two positive Fuchsian representation of $\Sigma_r.$ We will be interested in the representation:

\[
    \rho=(\rho_l,\rho_r):\pi_1(S)\to\text{Isom}_0(\A^{2,1})\simeq \text{PSL}(2,\R)\times\text{PSL}(2,\R).
\]

\begin{definition}
    Given a pair of positive Fuchsian representation $\rho_l,\rho_r:\pi_1(\Sigma_r)\to\text{PSL}(2,\R)$ we define $\Lambda(\rho)$ to be the graph in $\R\text{P}^1\times\R\text{P}^1$ of the unique $(\rho_l,\rho_r)-$equivariant orientation-preserving homeomorphism of $\R\text{P}^1,$ and $\Omega_\rho \coloneqq \Omega(\Lambda(\rho))$ its invisibile domain in $\A^{2,1}$
\end{definition}

Using the above construction, we can build examples of MGH spacetimes having holonomy any $\rho=(\rho_l,\rho_r)$ of this form. 

\begin{proposition}\label{examplesgenustwo}
    The domain $\Omega_\rho$ is invariant under the isometric action of $\pi_1(\Sigma_r)$ on $\A^{2,1}$ induced by $\rho.$ Moreover $\pi_1(\Sigma_r)$ acts freely and properly discontinuosly on $\Sigma_\rho$ and the quotient is a MGH spacetime of genus $r$ and holonomy $\rho.$ 
\end{proposition}
\begin{proof}
    By the definition of $\varphi$ and the action of $\PSL\times\PSL$ we have that $\Lambda_(\rho)$ is invariant under the action of $(\rho_l(\gamma),\rho_r(\gamma))$, for every $\gamma\in\pi_1(\Sigma_r).$ We recall from Corollary \ref{454} that $\Omega_\rho$ is the set of elements $x\in\PSL,$ such that $x\circ\varpi$ has no fixed point on $\S,$ the invariance of $\Omega_\rho$ also follows immediately, indeed: 
    \[
        (\rho_l(\gamma)\circ x\circ \rho_r(\gamma)^{-1})\circ \varphi = \rho_l(\gamma)\circ (x\circ\varphi)\circ\rho_2(\gamma)
    \]
    acts freely on $\S$ if $x\circ\varphi$ does. \\
    Let us show that for a compact set $K$ in $\Omega_\rho$, $\rho(\gamma)\cdot K$ stays in a compact region of $\Omega_{\rho}$ only for finitely many $\gamma\in\pi_1(\Sigma_r).$  This will also show that the action is free, since $\pi_1(\Sigma_r)$ has no torsion. For this purpose, take a sequence $x_n\in K$ and a sequence $\gamma_n\in \pi_1(\Sigma_r)$ not definitely constant. We claim that up to a subsequence, $(\rho(\gamma_n)\cdot x_n)$ converges to some $(\xi_+,\varphi(\xi_+))$ in $\Lambda(\rho).$ We will apply the criterion of convergence in $\partial\A^{2,1}$ seen in Lemma \ref{convergenza}.\\
    Since Fuchsian representations act cocompactly \textcolor{red}{means? and its true because?} on $\H^2$, the sequence $\rho_l(\gamma_n)$ has no converging subsequence in $\PSL$. By a know dynamical property of $\PSL$ (for more information view [BZ06]), up to taking a subsequence, there exist $\xi_-,\xi_+\in\S$ such that $\rho_l(\gamma_n)^{\pm 1}(\xi)\to\xi_{\pm}$ for all $\xi\neq \xi_{\mp}$ and that the convergence is uniform on compact sets of $(H^2\cup\S)\setminus\{\xi_{\mp}\}.$ By the equivariance condition, the same holds for $\rho_r(\gamma_n)$ where not $\xi_\pm$ are replaced by $\varphi(\xi_\pm).$\\
    To apply the criterion of Lemma \ref{convergenza}, pick any $p\in\H^2$, and recall taht $\rho(\gamma_n)\cdot x_n=\rho_l(\gamma_n)\circ x_n\circ \rho_r(\gamma_n)^{-1}.$ By the dynamical propery above, for any $\delta>0$ one can find $n_0$ such that $\rho_r(\gamma_{n})^{-1}(p)$ is in the $\delta$-neighborhood of $\varphi(\xi_-)$ (for the Euclidean metric on the closed disc), say $U_\delta$. Since $x_n$ lies in a compact region of $\Omega_\rho$, we can assume that is converges to $x_\infty \in \Omega_\rho$, hence $x_\infty\circ\varphi$ has no fixed point, and in particular $x_\infty\circ\varphi(\xi_-)\neq \xi_-.$\\
    Up to taking $\delta$ sufficiently small and $n_0$ large, $x_n(U_\delta)$ lies in a neighborhood $V_\epsilon$ of $x_\infty\circ\varphi(\xi_-)$ such that the closure of $V_\epsilon$ is disjoint from $\xi_-$. By construction $x_n\circ \rho_r(\gamma_n)^{-1}(p)$ converges to $\xi_+$ for $n$ large. The very same argument then shows that $(\rho(\gamma_n)\cdot x_n)^{-1}(p)=\rho_r(\gamma_n)\circ x_n^{-1}\circ\rho_l(\gamma_n)^{-1}(p)$ converges to $\varphi(\xi_+)$. And this concludes the claim. \\
    Finally, past and future boundary components $\partial_\pm\mathcal{C}(\Lambda(\rho))$ are contained in $\Omega_{\rho}$, since $\Lambda(\rho)$ is the graph of an otientation-preserving homeomorphism \textcolor{blue}{see remark 465}. Hence they are $\rho-$ invariant properly embedded Cauchy surfaces in $\Omega_\rho$ and project to Cauchy surfaces of the quotient of the action of $\rho(\pi_1(\Sigma_r)),$ which are homeomorphic to $\Sigma_r$. This shows that the quotient is a globally hyperbolic spacetime of genus $r$, which is maximal by Proposition \ref{513}.
\end{proof}

\subsection{Classification of genus $r\geq 2$.} We will now show that the examples of Proposition \ref{examplesgenustwo} are \textit{all} the MGH spacetimes of genus $r$. 

\begin{lemma}
    Let $\rho:(\rho_l,\rho_r)$ be a pair of Fuchsian representations, and $\varphi:\R\text{P}^1\to\R\text{P}^1$ be the unique $(\rho_l,\rho_r)-$equivariant orientation-preserving homeomorphism of $\R\text{P}^1.$ Then $\Lambda(\rho)$ is the unique proper achronal meridian in $\partial\A^{2,1}$ invariant under the action of $\pi_1(\Sigma_r)$ induced by $\rho.$
\end{lemma}

\begin{proof}
    Let $\Lambda$ be a proper achronal meridian invariant under the action of $\pi_1(\Sigma_r)$. We claim that the intersection $\Lambda\cap \Lambda_\varphi$ is not empty. Once the claim will be showed, the proof is concluded in the following way. If $(\xi_0 , \varphi(\xi_0))\in\Lambda,$ then: 
    \[
        (\rho_l(\gamma)\cdot\xi_0,\varphi(\rho_l(\gamma)\cdot\xi_0))=(\rho_l(\gamma)\cdot\xi_0,\rho_r(\gamma)\cdot\varphi(\xi_0))\in\Lambda.
    \] 
    However the $\rho_l(\pi_1(\Sigma_r))$-orbit of $\xi_0$ is dense in $\S,$ hence we deduce that $\Lambda$ contains $\Lambda_\varphi$. But both $\Lambda$ and $\Lambda_\varphi$ are homeomorphic to $\mathbb{S}^1$, which necessarily implies $\Lambda=\Lambda_\varphi.$\\
    Let us then show the claim. Let $\gamma$ be a non-trivial element in $\pi_1(\Sigma_r)$. It is knows that $\rho_l(\gamma)$ and $\rho_r(\gamma)$ are necessarily loxodromic elements in $\PSL$, hence we denote by $\xi_l^+(\gamma)$, and $\xi_r^+(\gamma)$ their attractive fixed points respectively. Notice that $\xi_r^+(\gamma)=\varphi(\xi_l^+(\gamma))$, hence: 
    \[
        (\xi_l^+(\gamma),\xi_r^+(\gamma))\in \Lambda_\varphi.
    \]
    By homological reasoning the curve $\Lambda$ must meet the leaf of the left ruling of $\partial\A^{2,1}$: 
    \[
        \lambda_{\xi_r^+(\gamma)}=\{(\eta,\xi_r^+(\gamma))\;|\;\eta\in\S\}.
    \]
    That is, there exists $\eta_0\in S$ such that $(\eta_0,\xi_r^+(\gamma))$ lies in $\Lambda$. But then $(\rho_l(\gamma))^k\eta_0,\xi_r^+(\gamma)$ lies in $\Lambda$ for $k>0$. If $\eta_0\neq \xi_l^-(\gamma)$ we can pass to the limit on $k$ and deduce that $(\xi_l^+(\gamma),\xi_r^+(\gamma))$ lies in $\Lambda.$ \\
    So far, the choice of $\gamma$ was arbitrary. To conclude, assume now by contradiction that for every $\gamma\in\pi_1(\sigma_r)$ the point $(\xi_l^-(\gamma),\xi_r^+(\gamma))$ lies in $\Lambda$. Take $\alpha,\beta\in\pi_1(S)$ so that the axes of $\rho_l(\alpha)$ and $\rho_l(\beta)$ do not intersect. We may assume that the cyclic order of end-points of those axes is 
    \begin{equation}\label{57}
        \xi_l^+(\alpha)<\xi_l^+(\beta)<\xi_l^-(\beta)<\xi_l^-(\alpha).
    \end{equation}
    Now given that $\xi_r^\pm(\alpha)=\varphi(\xi_l^\pm(\alpha))$ and the same hold for $\beta$ we have: 
\begin{equation}\label{58}
    \xi_r^+(\alpha)<\xi_r^+(\beta)<\xi_r^-(\beta)<\xi_r^-(\alpha).
\end{equation}
On the other hand, by assumption (applied to $\alpha,\beta$ and their inverses) the curve $\Lambda$ contains $(\xi_l^+(\alpha),\xi_r^-(\alpha)),(\xi_l^+(\beta),\xi _r^-(\beta)), (\xi_l^-(\beta ),\xi_r^+(\beta ))$ and $(\xi_l^-(\alpha),\xi_r^+(\alpha).$ By achronality of $\Lambda$, the cyclic order to the secondo comenents must be the same as that of the first components, hence from \refeq{57} we obtain: 
\[
    \xi_r^-(\alpha)<\xi_r^-(\beta)<\xi_r^+(\beta)<\xi_r^+(\alpha),
\]
which contradicts \refeq{58}.
\end{proof}

Given a pair $\rho=(\rho_l,\rho_r)$ of positive Fuchsian representations of $\pi_1(\Sigma_r),$ we denote by $M_\rho$ the MGH spacetime $\Omega_\rho/\rho(\pi_1(\Sigma_r))$ constructed in Proposition \ref{examplesgenustwo}.

\begin{corollary}
    For any pair $\rho=(\rho_l,\rho_r)$ of positive Fuchsian representations of $\pi_1(\Sigma_r)$, $M_\rho$ is the unique MGH spacetime with holonomy $\rho$.
\end{corollary}

We are only left with one last step for the classification result, we want to show that the left and right holonomies are necessarily positive Fuchsian. 

\begin{proposition}
    Let $M$ be an oriented, time-oriented, globally hyperbolic spacetime of genus $r\geq 2$ and let us endow a Cauchy surface $\Sigma$ with the orientation induced by the future normal vector. Then the left and right components of the holonomy $\rho=(\rho_l,\rho_r):\pi_1(\Sigma)\to\text{PSL}(2,\R)\times\text{PSL}(2,\R)$ are positive Fuchsian representation. 
\end{proposition}

\begin{observation}
We refer to the honolomy $\rho$ with respect to an orientation-preserving developing map. Therefore $\rho$ is well-defined up to conjugacy in $\text{PSL}(2,\R)\times\text{PSL}(2,\R)$
\end{observation}
\begin{proof}
    We will prove that the $\S$-flat bundles with holonomy $\rho_l$ and $\rho_r$ are isomorphic to the unit tangent bundle of $\Sigma.$ For the sake of definiteness, let us focus on $\rho_l$. We will construct an isomorphism: 
    \[
        \Phi_l:T^1\widetilde{\Sigma}\to \widetilde{\Sigma}\times\S
    \]
    equivariant with the respect to the action on $T^1\widetilde{\Sigma}$ by the actions by deck transformation, and the diagonal action given by $\rho_l$ on $\widetilde{\Sigma}\times\S.$ \\
    We define $\Phi_l$ as follows: for an element $(x,v)\in T^1\widetilde{\Sigma}$, let 
    \[
        \xi(x,v)=(\xi^l(x,v),\xi^r(x,v))\in \T
    \]
    be the end-point of the spacelike geodesic ray $\exp_x(tv)$ in $\A^{2,1}$ for positive $t$, Then we define $\Phi_l(x,v)=(x,\xi^l(x,v)).$ This map is clearly continuous, proper, equivariant and fiber preserving. \\
    In order to prove that it a bijection it suffices to notice that for any $x\in\widetilde{\Sigma}$ the map $\xi_x:T_x^1(\widetilde{\Sigma})\to\T$ is an embedding with image the boundary of the totally geodesic plane tangent to $\widetilde{\Sigma}$ at $x$. This boundary is the graph of an orientation-preserving map of $\S,$ so the projection $v\to\xi^l(x,v)$ is bijective. Moreover, by our choice of the orientation on $\Sigma$, the orientation on $T_x^1\widetilde{\Sigma}$ corresponds to the orientation induced on $\xi_x(T^1_x\widetilde{\Sigma})$ as graph of an orientation-preserving homeomorphism. The proof for $\rho_r$ is analogous.
\end{proof}


We can now state the classification result. We denote the \textit{deformation space} of MGH spacetimes of genus $r$ by:

\[
    \mathcal{MGH}(\Sigma_r)=\{g\;\text{MGH AdS metric on}\;\Sigma_r\times\R\}/\text{Diff}_0(\Sigma_r\times\R)
\]
where the group of diffeomorphism isotopic to the identity is acting by pull-back. The holonomy map takes value in the space of representation of $\pi_1(\Sigma_r)$ into $\PSL\times\PSL$ and is well-defined on the quotient $ \mathcal{MGH}(\Sigma_r).$
As a consequence of Proposition [...], the left and right components of the holonomy of elements of $ \mathcal{MGH}(\Sigma_r)$ are positive Fuchsian representations. The space of these representations up to conjugacy is idetified with the Teichm\"uller space of $\Sigma_r$ by the aforementioned work of Goldman []:
\[
    \mathcal{T}(\Sigma_r)\simeq\{\rho:\pi_1(\Sigma_r)\to\PSL\;\text{positive Fuchsian representation}\}/\PSL
\]

Therefore the holonomy map can be considered as a map from $ \mathcal{MGH}(\Sigma_r)$ with values in $\mathcal{T}(\Sigma_r)\times\mathcal{T}(\Sigma_r).$ Restating the classification in the original set of Mess \cite{Mess}: 
\begin{theorem}
    The holonomy map $\rho:\mathcal{MGH}(\Sigma_r)\to\mathcal{T}(\Sigma_r)\times\mathcal{T}(\Sigma_r)$ is a homeomorphism.
\end{theorem}
\section{Gauss map of spacelike surfaces}
    We will now introduce the \textit{Gauss Map} associated to a spacelike surface in Anti de Sitter space, this tool will be  useful in the study of the relation of Anti-de Sitter geometry and Teichm\"uller theory and hyperbolic surfaces.

    We will start by fixing notation and recalling some general facts about immersed spacelike surfaces in Anti de Sitter space. We will assume that all our immersed surfaces are of class $C^1$. \\
    Given a regular immersion $\sigma: S\to\A^{2,1}$ we will say that it is \textit{spacelike} if the pull-back metric $I\coloneqq\sigma^*g_{\A^{2,1}}$ is Riemannian, we call $I$ firs fundamental form of $\sigma$. 

The tangent bundle $TS$ is naturally identified to a subbundle of $T\A^{2,1}$ by means of $d\sigma$. The normal bundle $N_\sigma$ is then defined as the $g_{\A}$-orthogonal complement of $TS$ in $T\A^{2,1}$, The restriction of $g_{\A}$ to $N_\sigma$ is negative definite. Using the $g_\A$ induced decomposition: 
\[
\sigma^*T\A^{2,1}=TS\oplus N_\sigma,    
\]

the pull-back of the ambient Levi-Civita connection $\nabla$, restricted to sections tangent to $S$ splits as the sum of the Levi-Civita connection $\nabla^I$ of the first fundamental form $I$, and a symmetric 2-form with value in $N_\sigma$. As a consequence of time-orientability of $\A^{2,1}$, the normal bundle admits a natural trivialization, which is the same as a choice of a continuous unit normal vector field for $\sigma.$ We will denote by $\nu:S\to N_\sigma$ the future-directed unit normal section, and consider the decomposition for all vector field $X,Y$ tangent to $S$: 
\[
    \nabla_V W=\nabla^I_V W+II(V,W)\nu,
\]
where the symmetric $(2,0)$ tensor $II$ is called \textit{second fundamental form}. It will reveal useful to consider the $I-$symmetric $(1,1)$-tensor $B\in (TS)^*\otimes TS$ defined by $II(V,W)=I(B(V),W)$ which is called the \textit{shape operator} of $\sigma.$ As in the Riemannian case, it turns out that $\sigma_*(B(v))=\nabla_v\nu.$\\  
The first and second fundamental form of an immersion $\sigma$ satisfy constraint equations, known as \textit{Gauss-Codazzi equations}. More preciselly the Gauss equation consists of the identity: 
\begin{equation}\label{Gauss}
    K_I=-1-\det_I II
\end{equation}
where $K_I$ is the curvature of $I$ and $\det_I II$ is $\det B$ by definition. On the other hand the Codazzi equation states that $\nabla^III$ is a totally symmetric $3-$form. We have: 
\begin{equation} \nabla 
\end{equation}

\textcolor{red}{mancano eq. di codazzi}

What can be shown is that the Gauss-Codazzi equations relating first and second fundamental forms are \textit{necessary} but also \textit{sufficient}, more explicitly we mean:

\begin{theorem}\label{immcondition}
    Let $S$ e a simply connected surface, let $I$ be a Riemannian metric on $S$ and $II$  be a symmetric $(2,0)-$tensor on $S$. If $I$ and $II$ satisfy the Gauss-Codazzi equations \refeq{Gauss} and  \textcolor{red}{blablabla}
\end{theorem}

\subsection{Germs of spacelike immersions in Ads manifolds.}

Let us now consider the case of an oriented surface $\Sigma,$ not necessarily simply connected. Given a spacelike immersion $\sigma:\Sigma\to(M,g)$ where $(M,g)$ is an oriented Anti-de Sitter manifold, we can associate to $\sigma$ the pair $(I,II)$ as done in the previous section, where $II$ is computed with respect to the future unit normal vector $\nu$ of $\sigma$. We will always assume that the orientation of $\Sigma$ and $\nu$ are compatible with the orientation on $M.$\\
The pair $(I, II)$ only depends on the \textit{germ} of the immersion $\sigma.$\\
Given a pair $(I,II)$ on a surface $\Sigma$, we can perform the following construction. Let $\pi:\tilde{\Sigma}\to\Sigma$ be a universal cover, it follows that the pair $(\pi^*I,\pi^*II)$ satisfies the Gauss-Codazzi equations on $\tilde{\Sigma},$ hence by the existance part of Theorem \ref{immcondition} there exist a spacelike immersion $\tilde{\sigma}:\tilde{\Sigma}\to\A^{2,1}$ having \textit{immersion data} $(\pi^*I,\pi^*II)$. The uniqueness part of Theorem \ref{immcondition} has two main consequences: 

\begin{itemize}
    \item Any two such immersion differ by post-composition by a global isometry of $\A^{2,1}$. 
    \item Given any such $\tilde{\sigma},$ there exists a map $\rho:\pi_1(\Sigma)\to\text{Isom}_0(\A^{2,1})$ such that, for every $\gamma\in\pi_1(\Sigma),$ $f\circ\gamma=\rho(\gamma)\circ f$. 
\end{itemize}

More can be proved, $\rho$ is a group representation and changing $\tilde{\sigma}$ by post-composition with an isometry $f$ has the effect of conjugating $\rho$ by $f$. The immersion $\sigma$ can then be extended to an immersion of $U$, an open neighborhood of $\Sigma\times\{0\}$ in $\Sigma\times\R$, into $\A^{2,1}$, by mapping $(x,t)$ to the point $\gamma(t)$ on the timelike geodesic $\gamma$ such that $\gamma(0)=\sigma(p)$ and $\gamma^{\prime}(0)$ is the future normal vector of $\sigma$ at $x$. We want to explicit the expressions of the Anti-de Sitter metric in such a tubular neighborhood of $\sigma$ as it will be useful in the following: 

\begin{lemma}
    

\begin{equation}\label{64}
    4
\end{equation}


\end{lemma}

In short, given a pair $(I,II)$, the expression \refeq{64} provides a Lorentianz metric of constant curvature $-1$ on an open set $U$in $\Sigma\times\R$ containing the slice $\Sigma\times\{0\}$, and thus a germ of immersion of $\Sigma$ into an Anti-de Sitter three-manifold with immersion data $(I,II)$.\\
We want now to move on the case where $\Sigma$ is a closed surfaces. In such instance, the equivariant immersion $\tilde{\sigma}$ via the representation \textcolor{red}{scritto male qua} is necessarily an embedding, which can be extended to an embedding of $\tilde{\Sigma}\times\R$ onto a domain of dependence in $\A^{2,1}.$
The representation $\rho:\pi_1(\Sigma)\to\T$ coincides with the holonomy of a maximal globally hyperbolic Anti-de Sitter manifold $(M,g)$ (after the identification between $\pi_1(\Sigma)$ with $\pi_1(M)$ via embedding of $\Sigma\to M\simeq\Sigma\times\R)$, therefore $\rho$ consists of a pair of positive Fuchsian representation by Proposition \textcolor{red}{5.5.3}. \\ 

\subsection{Gauss map and projections} We can finally define our map for spacelike surfaces in $\A^{2,1}$. We have seen in \textcolor{red}{daje} that the space of timelike geodesic of $\PSL$ is naturally identified with $\H^2\times\H^2$, where the identification maps a geodesic of the form
\[
    L_p,q=\{X\in\PSL\;|\;X\cdot q =p\}
\]
to the pair $(p,q)\in \H^2\times\H^2$. We still suppose that the spacelike immersion is $C^1$ for the moment, we will deal later with the case of weaker regularity. 
\begin{definition}
    Let $\sigma:S\to\A^{2,1}$ be a spacelike immersion. The \textit{Gauss map} of $\sigma$ is the map $G_\sigma:S\to\H^2\times\H^2$ defined by $G_\sigma(x)=(p,q)$ such that $L_{p,q}$ is the timelike geodesic orthogonal to $\text{Im}(d_x\sigma))$ at $\sigma(x).$\\
    It follows from the equivariance shown in \textcolor{red}{inserire}, that the Gauss map $G_\sigma$ is natural with respect to the action of the isometry group, namely: 
    \[
        G_{f\circ\sigma}=f\cdot G_\sigma
    \]
    for every $f\in\text{Isom}_0(\A^{2,1})=\PSL\times\PSL$.
    \end{definition}

    \begin{example}
        In Lemma \textcolor{red}{forse va messo qua e basa} we gave an isometric embedding of $\H^2$ in $\A^{2,1}$ with image the plane $P_{\text{Id}}$ dual to the identity. This isometric embedding is defined by sending $p\in\H^2$ to the unique
    \end{example}

    By construction, the Gauss map of a spacelike immersion $\sigma$ is invariant by reparametrization, in the sense that $G_{\sigma\circ\phi}=G_\sigma\circ\phi$ for a diffeomorphism $\phi:S^{\prime} \to S$. Hence it makes sense to talk about \textit{the} Gauss map of a spacelike surface in $\A^{2,1}$. For example, for the plane dual to the identity. 
    
    \begin{lemma}\label{633}
        Given a spacelike immersion $\sigma:S\to\A^{2,1}$, with future unit normal vector field $\nu$, if $\sigma(p)=\text{Id}$, then 
        \begin{equation}\label{gaussmap}
            G_\sigma(p)=G_{P_\text{Id}}(\exp(\frac{\pi}{2}\nu(p)))
        \end{equation}
    \end{lemma}
    \begin{proof}
        It's a consequence of the example. We need to observe that the geodesic leaving from Id with velocity $\nu(p)$ meets orthogonally $P_\text{Id}$ at $\exp((\pi/2)\nu(p))$
    \end{proof}

    We want to consider now the map: 
    \[
        \text{Fix}:T_{\text{Id}}^{1,+}\A^{2.1}\to\H^2
    \]
    where $\hf$ denotes the hyperboloid of future unit timelike vectors in $T_{\text{Id}}\A^{2,1}$, such that $\text{Fix}(\nu)$ us the fixed point of the one-parameter elliptic group $\{\exp(t\nu)\;|\;t\in\R\}$. This map is equivariant for the action of $\PSL$, which acts on the hyperboloid $\hf$ by the adjoint representation and on $\H^2$ by the obvious action. Since both $\hf$ and $\H^2$ have constant curvature $-1$, it follows from equivariancy that $\text{Fix}$ is an isometry.  \\
    In terms of the map $\textit{Fix},$ Equations \refeq{gaussmap} read as: 

    \begin{equation}\label{66}
        G_\sigma(p)=(\text{Fix}(\nu(p)),\text{Fix}(\nu(p))),
    \end{equation}

    provided that $\sigma(p)=\text{Id}.$\\ Via Lemma \ref{633} and the naturality, we can recover a different description of the Gauss map. 

    \begin{lemma}
        Given a spacelike immersion $\sigma: S\to \A^{2.1}$ with future unit normal vector field $\nu$,
        \[
            G_\sigma(p)=(\text{Fix}((R_{\sigma(p)^{-1}})_* (\nu(p)),\text{Fix}((L_{\sigma(p)^{-1}})_* (\nu(p))).
     \]
    \end{lemma}

    \begin{proof}
        If $\sigma(p)=\text{Id}$, the the equality holds true by Equation \refeq{66}. For the general case, the immersion $\sigma^{\prime} =(\text{Id},\sigma(p))\circ\sigma$ has the property that $\sigma^{\prime}(p)=\text{Id}$, and the future normal vector at $\sigma^{\prime} (p)$ equals $\nu^{\prime} (p)=(R_{\sigma(p)^{-1}}))_{\ast} (\nu(p))$. Therefore: 
        \[
            G_{\sigma^{\prime} }(p)=(\text{Fix}((R_{\sigma(p)^{-1}})_{\ast} (\nu(p)),\text{Fix}((R_{\sigma(p)^{-1}})_{\ast} (\nu(p))).
        \]

        Now by the naturality of the Gauss map it follows, 
        
        \begin{align*}
            G_\sigma(p)&=(\text{Id},\sigma(p)^{-1})\cdot G_{\sigma^{\prime}}(p) \\
            &=(\text{Fix}((R_{\sigma(p)^{-1}})_{\ast} (\nu(p)),\sigma(p)^{-1}\circ\text{Fix}((R_{\sigma(p)^{-1}})_{\ast} (\nu(p))) \\
            &=(\text{Fix}((R_{\sigma(p)^{-1}})_{\ast} (\nu(p)),\text{Fix}((L_{\sigma(p)^{-1}})_{\ast} (\nu(p))).
        \end{align*}
        
        where in the last line we used that Fix is equivariant with respect to the adjoint action on the hyperboloid $\hf$.
        
    \end{proof}

    We will refer to the components of the Gauss map as the \textit{left} and \textit{right projections,} and will denote them accordingly: $\Pi_{l},\Pi_{r}:S\to\H^2.$\\

\begin{observation}
    Under the identification given by Fix, the left and right projections can be interpreted in the following way. Given $p\in S$, $\Pi_l(p)$ is the parallel transport in Id of the future unit vector $\nu(p)$ at $\sigma(p)$ with respect to the right-invariant connection $D^r$. The right projection is then obtained using the left-invariant connection. 

    \textcolor{red}{In realtà Diaf fa il remark sotto 6.3.6}
\end{observation}


     \begin{observation}
        In the work of Mess, we can find yet a different interpretation of the Gauss map. Given $p\in S$, one can find a unique left isometry $f_l(p)$ and a unique right isometry $f_r(p),$ sending the tangent plane $P$ to the image of $\sigma$ at $\sigma(p)$ to $P_{\text{Id}}$. Indeed the isometries $f_l(p)$ and $f_r(p)$ are siply obtained by left and right multiplicationby the inverse of dual point of the tangent plane $P$, namely $\zeta(p)=\exp_{\sigma(p)}((\pi/2)\nu(p))$. Using the identification of the dual plane $P_{\text{Id}}$ with $\H^2$, provided by Lemma \textcolor{red}{uk}, $\Pi_l(p)$ and $\Pi_r(p)$ are the image of $\sigma(p)$ under the right and left isometries, respectively:
        \[
        \Pi_l(p)=f_r(p)\circ\sigma(p)=(\text{Id},\zeta(p))\cdot \sigma(p)\;\text{and}\;\Pi_r(p))=f_l\circ\sigma(p)=(\zeta(p)^{-1},\text{Id})\cdot\sigma(p)\
        \]
    \end{observation}


We want now to prove formulae which express the pull-back of the hyperbolic metrics by the left and right projections. When applying these formulae to the embedding data of a surface in an MGH Cauchy compact Anti-de Sitter spacetime $(M,g)$, we obtain a pair of hyperbolic metrics whose isotopy classes are the parameters of $(M,g)$
in $\mathcal{T}(S)\times\mathcal{T}(S).$

\begin{proposition}
    Let $\sigma:S\to\A^{2,1}$ be a spacelike immersion, let $\Pi_l,\Pi_r:S\to\H^2$ be the left and right projections and let $g_{\H^2}$ be the hyperbolic metric. Then 
    \begin{equation}
        \Pi_l^*
    \end{equation}
\end{proposition}

\begin{proof}
    
\end{proof}

\textit{Non-smooth surfaces.} The construction of the Gauss map can be extended (\textit{how?}) in the non-smooth setting, for instance for convex spacelike surfaces $S\in\A^{2,1}$. Then one define the set-valued Gauss map as the map sending each $x\in S$ to the set of future unit vectors in $T_x^{1,+}\A^{2,1}$ orthogonal to support planes of $S$ at $x$. Hence the image of $x$ is a convex subset of $T_x\A^{2,1}$, and it reduces to a single point if and only if $S$ is differentiable at $x$. The image of $G$ in $T^{1,+}_x\A^{2,1}$ is a $C^{1.1}$ surface. \\


%\include{conclusione}
%\appendix
%\include{appendice}
\printbibliography
\end{document}
